% concatenated from arxiv1015.tex
\documentclass[12pt]{article}
\usepackage{amsmath,amsfonts,latexsym,mathrsfs, amsthm,amssymb}
\pagenumbering{arabic}
\makeatletter
  \renewcommand{\theequation}{%
      \thesection.\arabic{equation}}
  \@addtoreset{equation}{section}
\makeatother
\newtheorem{proposition}{Proposition}[section]
\newtheorem{theorem}{Theorem}[section]
\newtheorem{lemma}{Lemma}[section]
\newtheorem{remark}{Remark}[section]
\newtheorem{corollary}{Corollary}[section]
\newtheorem{definition}{Definition}[section]
\newtheorem{example}{Example}[section]
\usepackage{color}
\usepackage{ulem}

%\textheight=24cm
%\voffset=-2.5cm
%\textwidth=17cm
%\hoffset=-1.5cm

\date{ }

\begin{document}

\title{Stochastic optimal transport for the Langevin dynamics and its zero--mass limit
}

\author{
Toshio Mikami\thanks{Partially supported by JSPS KAKENHI Grant Number 24K06765.
}\\
Department of Mathematics, Tsuda University
}


\maketitle

\begin{abstract}
We introduce a stochastic optimal transport for the Langevin dynamics with positive mass and study its zero--mass limit.
The new aspect of this paper is that we only fix the initial and terminal probability distributions of the positions of particles under consideration, but not those of their velocities with Heisenberg's uncertainty principle in mind.
In the zero--mass limit, we show that the minimizer of our stochastic optimal transport is tight if and only if the initial momentum (=(mass)$\times$(velocity)) of a particle converges to zero.
We also show that the limit of a minimizer of our stochastic optimal transport is a minimizer of a standard stochastic optimal transport for continuous semimartingales.
\end{abstract} 

Keywords:  stochastic optimal transport, Langevin equation, zero--mass limit

AMS subject classifications:  93E20, 82C31, 49Q22



\section{Introduction}\label{sec:1}

The optimal mass transport theory plays a crucial role in many fields
e.g.,  information sciences and metric measure space
(see e.g., \cite{1-PC, 1-RR,1-V2} and the references therein).

For $d\ge 1$, let $\mathcal{P}(\mathbb{R}^d)$ denote the space of all Borel probability measures on $\mathbb{R}^d$ endowed with weak topology, and let
$$\mu_1(dx):=\mu(dx\times \mathbb{R}^d),\quad\mu_2(dy):=\mu(\mathbb{R}^d\times dy),\quad
\mu\in \mathcal{P}(\mathbb{R}^d\times \mathbb{R}^d),$$
$$\Pi(P,Q):=\{\mu\in \mathcal{P}(\mathbb{R}^d\times \mathbb{R}^d):
\mu_1=P,\mu_2=Q\},\quad P,Q\in \mathcal{P}(\mathbb{R}^d).$$
Let $AC([0,1];\mathbb{R}^d)$ and $P^{X}$ denote the space of all absolutely continuous functions from $[0,1]$ to $\mathbb{R}^d$ and the probability distribution of a random variable $X$ defined on a probability space, respectively.
In this paper, the probability space under consideration is not fixed.
We use the same notation $P$ for different probabilities when it is not confusing.

For $P_0,P_1\in \mathcal{P}(\mathbb{R}^d)$, the following is a typical example of the Monge--Kantorovich problem
that is a class of optimal mass transports:
\begin{equation}
T(P_0,P_1):=\inf\left\{E[|Y-X|^2]:P^X=P_0, P^Y=P_1\right\}.
\end{equation}
The following holds:
\begin{eqnarray}T(P_0,P_1)&=&\inf\left\{E\left[\int_0^1\left|\frac{d}{dt}X(t)\right|^2dt\right]:X(\cdot)\in AC([0,1];\mathbb{R}^d), {\rm a.s. }, \right.\label{0.1}\\
&&\qquad \left.P^{(X(0), X(1))}\in \Pi(P_0,P_1)\right\}.\nonumber
\end{eqnarray}
%where $AC([0,1];\mathbb{R}^d)$ and $P^{X}$ denote the space of absolutely continuous functions from $[0,1]$ to $\mathbb{R}^d$ and  the probability distribution of a random variable $X$, respectively.
Indeed, for $\varphi (\cdot)\in AC([0,1];\mathbb{R}^d)$, by Jensen's inequality,
$$\int_0^1\left|\frac{d}{dt}\varphi (t)\right|^2dt\ge |\varphi (1)-\varphi (0)|^2,$$
where the equality holds if and only if 
$$\frac{d}{dt}\varphi (t)=\varphi (1)-\varphi (0),\quad dt{\rm -a.e.}$$ 
(see \cite{M09, M2021, 1-RR, 1-V2} and references therein).

If $T(P_0,P_1)$ is finite, then a minimizer of (\ref{0.1}) exists and satisfies the following:
\begin{equation}\label{1.2.0}
X(t)=X(0)+t\{X(1)-X(0)\},\quad 0\le t\le 1, \quad {\rm a.s.}.
\end{equation}%  \quad \frac{d^2}{dt^2}X(t)=0,$$
If, in addition, $P_0(dx)\ll dx$, then
there exists a convex function $\varphi$ on $\mathbb{R}^d$ such that 
\begin{equation}\label{1.3.0810}
X(1)=D\varphi(X(0)),\quad {\rm a.s.},
\end{equation}
where $D:=(\partial /\partial x_i)_{i=1}^d$.
In particular, $\{X(t)\}_{0\le t\le 1}$ is a measurable function of $(t,X(0))$
(see \cite{Bre1, Bre2,GangboMcCann,1-RR, 1-V2} and references therein).
%In particular, the following holds:
%\begin{eqnarray}\label{1.2.0}
%T(P_0,P_1)&=&\inf\left\{E\left[\int_0^1\left\{\left|\frac{d}{dt}X(t)\right|^2+\left|\frac{d^2}{dt^2}X(t)\right|^2\right\}dt\right]:\right.\\
%&&\qquad \left.\frac{d}{dt}X(\cdot)\in AC([0,1];\mathbb{R}^d), {\rm a.s. }, P^{(X(0), X(1))}\in \Pi(P_0,P_1)\right\}\nonumber
%\end{eqnarray}
%(see the Appendix for proof).


The so-called Schr\"odinger's problem in the theory of stochastic processes 
was discussed in E. Schr\"odinger's papers \cite{S1, S2}.
It is the origin of stochastic optimal transport (SOT for short) (see \cite{M2021} and the references therein) and 
is also considered an entropic regularized optimal mass transport  in data science
nowadays (see \cite{1-PC} and the references therein).

Let
$$g(t,x):=\frac{1}{\sqrt{2\pi t}^d}\exp \left( -\frac{|x|^2}{2t}\right),\quad (t, x)\in (0,\infty )\times \mathbb{R}^d.$$
For Borel probability measures $\mu,\nu$ on a topological space $S$, let 
$$H(\mu\|\nu):=
\begin{cases}
\displaystyle \int_S \log\frac{d\mu}{d\nu}(x)d\mu (x),&\mu\ll\nu,\\
\infty, &\hbox{otherwise.}
\end{cases}
$$

For $P_0,P_1\in \mathcal{P}(\mathbb{R}^d)$, the following is a typical example of Schr\"odinger's problem:
\begin{eqnarray}\label{1.3.-1}
&\inf\left\{H(\mu(dx\hbox{ }dy)\|P_0(dx)g(1,y-x)dy):\mu\in \Pi(P_0,P_1)\right\}.
\end{eqnarray}
If there exists $\mu\in \Pi(P_0,P_1)$ and a Borel measurable function $q^\mu$ on  
$\mathbb{R}^d\times \mathbb{R}^d$ such that 
$$\mu (dx\hbox{ }dy)=q^\mu(x,y)P_0(dx)g(1,y-x)dy,$$
then $P_1(dy)$ has a density
$$p_1(y):=\int_{\mathbb{R}^d}q^\mu(x,y)P_0(dx)g(1,y-x).$$
In particular, (\ref{1.3.-1}) is equal to the following (see e.g., \cite{1-PC}):
\begin{eqnarray}
&&\inf\left\{E\left[\frac{|Y-X|^2}{2}\right]+H(P^{(X,Y)}\|P^{X}\times P^{Y}):
P^X=P_0, P^Y=P_1\right\}\label{1.6.0}\\
&&\qquad +\frac{d}{2}\log (2\pi )+\int_{\mathbb{R}^d}p_1(y)\log p_1(y)dy,\nonumber
\end{eqnarray} 
provided it is well defined.
The infimum in (\ref{1.6.0}) is called an entropic regularized optimal mass transport  in data science.
In information science, $H(P^{(X,Y)}\|P^{X}\times P^{Y})$
is called the mutual information of $X$ and $Y$ and is denoted by $I(X;Y)$ (see e.g., \cite{Cover}).

(\ref{1.3.-1}) is also equal to the following, which is an  SOT  analog of (\ref{0.1}):
\begin{eqnarray}
&&\inf\left\{H(P^{X(\cdot)}\|P^{X(0)+W(\cdot)}):dX(t)=u_X(t)dt +dW(t), \right.\nonumber\\
&&\qquad \left. P^{(X(0), X(1))}\in \Pi(P_0,P_1)\right\}\nonumber\\
&=&\inf\left\{E\left[\int_0^1\frac{1}{2}\left|u_X(t)\right|^2dt\right]:dX(t)=u_X(t)dt +dW(t), \right.\label{1.3.0}\\
&&\qquad \left. P^{(X(0), X(1))}\in \Pi(P_0,P_1)\right\},\nonumber
\end{eqnarray}
where $u_X(t)$ and $W(t)$ are, respectively, an $\mathbb{R}^d$--valued progressively measurable stochastic process
and an $\mathbb{R}^d$--valued  Brownian motion defined on the same filtered probability space
(see \cite{D91, F88}, \cite{J75}--\cite{MT06}, \cite{Z86}--\cite{Z86-3}
%\cite{J75, L2, M04, M08, M09, M21, M2021, MT06, Z86, Z86-2, Z86-3} 
and the references therein).
 In this paper, we use the same notation $W$ for different  Brownian motions when it is not confusing.

Indeed, (\ref{1.3.-1}) is less than or equal to (\ref{1.3.0}) since if 
$$dX(t)=u_X(t)dt +dW(t),\quad dP^{X(\cdot)}\ll dP^{X(0)+W(\cdot)},$$
then 
\begin{eqnarray}
&&P^{(X(0),X(1))}(dxdy)\label{1.8.0}\\
&=&E\left[\frac{dP^{X(\cdot)}}{dP^{X(0)+W(\cdot)}}\biggl|(X(0), X(0)+W(1))=(x,y)\right]
P^{(X(0), X(0)+W(1))}(dxdy),\nonumber
\end{eqnarray}
and by Jensen's inequality,
\begin{eqnarray*}
&&H(P^{X(\cdot)}\|P^{X(0)+W(\cdot)})\\
&=&\int_{C([0,1];\mathbb{R}^d)}\left\{\frac{dP^{X(\cdot)}}{dP^{X(0)+W(\cdot)}} (\omega)\log\frac{dP^{X(\cdot)}}{dP^{X(0)+W(\cdot)}} (\omega)\right\}dP^{X(0)+W(\cdot)}(\omega)\\
&\ge& E\left[E\left[\frac{dP^{X(\cdot)}}{dP^{X(0)+W(\cdot)}}\biggl|(X(0), X(0)+W(1))\right]\right.\\
&&\qquad \times \left.\log \left(E\left[\frac{dP^{X(\cdot)}}{dP^{X(0)+W(\cdot)}}\biggl|(X(0), X(0)+W(1))\right]\right)\right]\\
&=&H(P^{(X(0),X(1))}\|P^{(X(0), X(0)+W(1))}).
\end{eqnarray*}
There exist functions $u,v$ on $\mathbb{R}^d$ such that 
$u(\cdot)+|\cdot|^2/2$ and $v(\cdot)+|\cdot|^2/2$ are convex on $\mathbb{R}^d$
and that
%$$\mu_o (dx\hbox{ }dy)=\exp (u(x)+v(y))\frac{1}{\sqrt{2\pi}^d}\exp \left( -\langle x,y\rangle\right)P_0(dx)p_1(y)dy$$
$$\mu_o (dx\hbox{ }dy):=\exp (-u(x)-v(y))g(1,y-x)P_0(dx)p_1(y)dy\in \Pi(P_0,P_1)$$
(see \cite{1-Jamison1974}).
The following also has a unique weak solution:
\begin{eqnarray*}
dX_o(t)&=&D_x \log \left(\int_{\mathbb{R}^d}g(1-t, y-X_o(t))\exp (-v(y))p_1(y)dy\right)dt +dW(t),\quad 0<t<1,\\
P^{X_o(0)}&=&P_0,
\end{eqnarray*}
and the following holds:
$$dP^{X_o(\cdot)}(\omega)=\exp (-u(\omega(0))-v(\omega(0)+\omega(1)))p_1(\omega(0)+\omega(1)) dP^{X_o(0)+W(\cdot)}(\omega)$$
(see \cite{J75}).
In particular, from (\ref{1.8.0}),
\begin{eqnarray*}
P^{(X_o(0),X_o(1))}&=&\mu_o,\\
H(P^{X_o(\cdot)}\|P^{X_o(0)+W(\cdot)})
&=&H(P^{(X_o(0),X_o(1))}\|P^{(X_o(0), X_o(0)+W(1))})\\
&=&H(\mu_o\|P_0(dx)g(1,y-x)dy).
\end{eqnarray*}
If (\ref{1.3.-1}) is finite, then $\mu_o$ is  a unique minimizer of (\ref{1.3.-1})
(see \cite{1-Jamison1974, 1-RT} and also \cite{M2021}),
which implies that (\ref{1.3.-1}) is greater than or equal to (\ref{1.3.0}) and $X_o$ is a unique minimizer of (\ref{1.3.0})
(see \cite{D91, F88, J75},
\cite{M08}--\cite{MT06}, \cite{Z86}--\cite{Z86-3} and the references therein).

\begin{remark}\label{rk1.1.0}
In \cite{M04}, we showed that if $P_0(dx)\ll dx$, and $P_0$ and $P_1$ have the second moments, then 
the zero--noise limit of $X_o$ exists, and the limit is a gradient of a convex function that satisfies 
(\ref{1.3.0810}).
It is a probabilistic proof of Monge's problem
(see \cite{Bre1,Bre2,L2,M09,M2021,1-MTsiam,1-PC,1-RR,1-V2} and the references therein
and also Remark \ref{rk1.3} given later).

\end{remark}








The Langevin equation is a generalization of the Newton equation of motion and describes the motion of a particle subject to friction and stochastic forcing.

Let  $k_B$ and $T$ denote Boltzmann's constant and the absolute temperature, respectively.
For $m, \gamma >0$ and a sufficiently smooth $U:\mathbb{R}^d\rightarrow\mathbb{R}$,
the following is a class of Langevin equations with positive mass: for $t\in (0,1)$,
\begin{equation}\label{0}
%dX(t)&=&\frac{d}{dt}X(t)dt,\\
m\cdot d\left(\frac{d}{dt}X(t)\right)=\left\{-DU(X(t))-\gamma \left(\frac{d}{dt}X(t)\right)\right\}dt+\sqrt{2\gamma k_BT}dW(t)
\end{equation}
(see e.g., \cite{IW14} for the SDE).
Here $m$, $U$, and $\gamma$ denote 
the mass of a particle, an interaction potential function, and
the friction coefficient, respectively.
$\gamma dX(t)/dt$ and $\sqrt{2\gamma k_BT}dW(t)/dt$ are also 
a linear dissipation and a stochastic forcing, respectively
(see e.g., \cite{PAV}).
If $\gamma=0$, then (\ref{0}) is a class of Newton equations of motion.

In (\ref{0.1}) and (\ref{1.3.0}), $X(\cdot)\in AC([0,1];\mathbb{R}^d)$ and 
$X(\cdot)$ is a semimartingale, respectively.
Inspired by the Langevin equation, we are interested in studying  an SOT in the case where 
$X(\cdot)\in AC([0,1];\mathbb{R}^d)$ and  %\quad \frac{d}{dt}X(t)\not\in AC([0,1];\mathbb{R}^d),$$
 $dX(t)/dt$ is a semimartingale.
 
In this paper, we introduce an SOT for a class of Langevin dynamics with positive mass $m$ and study its zero--mass limit, i.e., the limit as $m\to 0$.
We consider the case where the initial and terminal probability distributions of the positions of  particles under consideration are fixed
and where the initial momentum converges to zero as $m\to 0$.
We emphasize that the initial velocity does not necessarily converge to zero as $m\to 0$.
Since we are interested in the zero--mass limit and since we assume that we know the initial and terminal probability distributions of the positions of particles,
we consider, with Heisenberg's uncertainty principle in mind, the case where we do not know the probability distributions of the momenta of particles.
Since the Langevin equation with positive mass is an SDE for the position and velocity of a particle, our problem is a new class of SOTs.
We refer readers to \cite{Birrell, FW-1, FW-2, Ishii} and the references therein for the zero--mass limit of the Langevin equation with a variable friction coefficient.

%Here, $P^X$ denotes the probability distribution of a random variable $X$. 


We describe our problem more precisely.
For notational simplicity, we consider the SDE for the position and momentum of a particle, instead of its velocity.
%Let $\mathcal{P}(\mathbb{R}^d)$ denote the space of all Borel probability measures on $\mathbb{R}^d$ endowed with weak topology.
Let $\sigma:[0,1]\times \mathbb{R}^{2d}\rightarrow M(d,\mathbb{R})$
 be a bounded Borel measurable $d\times d$--matrix function.
For $m>0$, let $\mathcal{A}^m$ denote the set of $\mathbb{R}^{d}\times \mathbb{R}^{d}$--valued continuous semimartingales $\{Z(t)=(X(t),Y(t))\}_{0\le t\le 1}$ defined on a complete filtered probability space such that the following holds weakly:
\begin{eqnarray}
dX(t)&=&\frac{1}{m}Y(t)dt,\label{1.1}\\
dY(t)&=&\left\{u_X(t)-\frac{\gamma }{m}Y(t)\right\}dt+\sigma(t,Z(t))dW(t),\quad 0< t<1.\label{1.2}%\\
%P^{X(t)}&=&P_t, \quad t=0,1.
\end{eqnarray}
Here  $\{u_X (t)\}_{0\le t\le 1}$ and $\{W(t)\}_{0\le t\le 1}$ are
a progressively measurable $\mathbb{R}^d$--valued stochastic process
and an $\mathbb{R}^d$--valued Brownian motion, respectively, defined on the same filtered probability space (see e.g., \cite{IW14}).
% If there exists $b:[0,1]\times \mathbb{R}^d\rightarrow\mathbb{R}^d$ such that
%$u_X(\cdot)=b(\cdot,X(\cdot))$, 
%then we write $b=b_{X}$.
% In this paper, the probability space under consideration is not fixed.
%When it is not confusing, we use the same notation $P$ and $W$ for different probabilities and Brownian motions, respectively.
We omit the dependence of $Z\in \mathcal{A}^m$ on $m$ when it is not confusing.


\begin{remark}\label{rk1.1}
(i) For $m>0$, (\ref{1.1})--(\ref{1.2}) is equivalent to (\ref{1.1}) and the following:
\begin{eqnarray}\label{1.6}
%dX(t)&=&\frac{d}{dt}X(t)dt,\\
m\cdot d\left(\frac{d}{dt}X(t)\right)&=&\left\{u_X(t)-\gamma \left(\frac{d}{dt}X(t)\right)\right\}dt\\
&&\qquad +\sigma\left(t,\left(X(t), m\frac{d}{dt}X(t)\right)\right)dW(t).\nonumber
\end{eqnarray}
(ii) For $\{Z(t)=(X(t),Y(t))\}_{0\le t\le 1}\in \mathcal{A}^m$,
$Y(t)$ is $\mathcal{F}^X_t:=\sigma [X(s);0\le s\le t ]$--measurable for $t\in (0,1]$, since
%such that 
%$E[\int_0^1 |u_X(t)|dt]$ is finite,
%$u_X(t)$ and $W(t)$ can be taken to be $\sigma [Y(0), X(s);0\le s\le t ]$--progressively measurable
%and to be an $\mathcal{F}^X_t:=\sigma [Z(s);0\le s\le t ]$--Brownian motion, respectively.
%Indeed, in (\ref{1.1})--(\ref{1.2}), 
%$$Y(t)-Y(0)-\int_0^t \left\{E[u_X(s)|\sigma [Y(0), X(\alpha);0\le \alpha\le s ]]-\frac{\gamma }{m}Y(s)\right\}ds
%$$
%is $\sigma [Y(0), X(s);0\le s\le t ]$--measurable for $t\in [0,1]$, since 
%$t\mapsto Y(t)$ is continuous, a.s., and 
$$Y(t)=
m \times \lim_{h\downarrow 0}\frac{X(t)-X(t-h)}{h},\quad 0<t\le 1, \quad {\rm a.s..}
$$
In the same way, $Y(0)$ is $\mathcal{F}^X_{0+}:=\cap_{t>0}\mathcal{F}^X_t$--measurable.
%In (\ref{1.2}), we use  a notation $u_X$, since 
%we only have to consider the case where $u_X(t)$ is $\sigma[Z(s); 0\le s\le t]$--measurable when we study our problems (\ref{1.3}) and (\ref{1.12}) given later,
%under (A1, iv) in section \ref{sec:2} (see e.g., \cite{M2021}, p. 41, Lemma 2.1).
%$$u_X(t)dt=
%dY(t)+\frac{\gamma }{m}Y(t)dt-\sigma(t,(X(t),Y(t)))dW(t)
%$$
%(see also Lemma \ref{lm3.2}).
\end{remark}



%We give notations to describe our SOT. 
For $m>0$, $B\subset \mathcal{P}(\mathbb{R}^d)$, and $P,Q\in \mathcal{P}(\mathbb{R}^d)$, let
%$$\mu_1(dx):=\mu(dx\times \mathbb{R}^d),\quad\mu_2(dx):=\mu(\mathbb{R}^d\times dx),\quad
%\mu\in \mathcal{P}(\mathbb{R}^d\times \mathbb{R}^d),$$
%$$\Pi(P,Q):=\{\mu\in \mathcal{P}(\mathbb{R}^d\times \mathbb{R}^d):
%\mu_1=P,\mu_2=Q\},\quad P,Q\in \mathcal{P}(\mathbb{R}^d),$$
\begin{eqnarray*}
\mathcal{A}^m(B,P;Q)&:=&\{\{ Z(t)=(X(t),Y(t))\}_{0\le t\le 1}\in \mathcal{A}^m:
P^{(X(0), X(1))}\in \Pi(P,Q), P^{Y(0)}\in B\},\\
\mathcal{A}^m(P_0;P_1)&:=&\mathcal{A}^m(\mathcal{P}(\mathbb{R}^d),P_0;P_1).
\end{eqnarray*}
%We also write
%$$\mathcal{A}^m_0(P;Q):=\mathcal{A}^m(\{\delta_0\},P;Q),$$
Let 
$$L:[0,1]\times \mathbb{R}^{2d} \times \mathbb{R}^d\longrightarrow [0,\infty )$$
be  Borel measurable.  
We also write
$$L(t,x,y;u):=L(t,z;u), \quad 
(t,u)\in [0,1]\times \mathbb{R}^d, z=(x,y)\in \mathbb{R}^d \times\mathbb{R}^d.$$

The following is an SOT for the Langevin dynamics when mass $m>0$.
\begin{definition}[SOT for the Langevin dynamics with positive mass\label{def1}]
For $m>0$, $B\subset \mathcal{P}(\mathbb{R}^d)$, and $P,Q\in \mathcal{P}(\mathbb{R}^d)$, let
\begin{eqnarray}
V^m(B, P_0;P_1)
&:=&\inf \left\{E\biggl[\int_0^1 L(t,Z(t);u_X (t))dt \biggr]:Z\in \mathcal{A}^m(B,P_0;P_1)\right\},\nonumber\\
\label{1.3}\\
V^m(P_0,P_1)&:=&V^m(\mathcal{P}(\mathbb{R}^d), P_0;P_1).\nonumber
\end{eqnarray}
%When $L=|u|^r, r>0$, we write $V=V_r$.
If the set over which the infimum is taken is empty, we set the infimum to be infinite.
\end{definition}
 
 
If $0<m_n\to 0, n\to\infty$, then any sequence $\{Z_n=(X_n, Y_n)\in \mathcal{A}^{m_n}(P_0,P_1)\}_{n\ge 1}$
such that the initial momentum $Y_n(0)=m_n dX_n(t)/dt|_{t=0}$ does not converges to $0$ as $n\to\infty$ is not tight since
the following is not continuous in $t$:
\begin{equation}\label{1.14}
\lim_{n\to \infty}\exp\left(-\frac{\gamma t}{m_n}\right)
=
\begin{cases}
0,&t\in (0,1],\\
1, &t=0
\end{cases}
\end{equation}
(see (\ref{3.9}), (\ref{3.11}), and also (\ref{3.7}) for notation).
This is one of the reasons  we restrict a class of $Y(0)$ under consideration, when we study the zero--mass limit (see section \ref{sec:2} for more discussion).


\begin{remark}\label{rk1.3}
% In (\ref{0}), if $\sqrt{2\gamma k_BT}=0$, then
%$$m\frac{d^2}{dt^2}X(t)=-DU(X(t))-\gamma \frac{d}{dt}X(t).$$
In (\ref{1.1})--(\ref{1.2}), if $\sigma(\cdot,\cdot)=0$,  then
$dX(\cdot)/dt\in AC([0,1];\mathbb{R}^d)$ and 
$$%m\frac{d}{dt}X(t)=Y(t),\quad 
m\frac{d^2}{dt^2}X(t)=u_X(t)-\gamma \frac{d}{dt}X(t).$$
In particular,  if %$\gamma=0$ and 
$\sigma(\cdot,\cdot)=0$, then 
\begin{eqnarray}
\quad V^m(P_0,P_1)
&=&\inf \left\{E\biggl[\int_0^1 L\left(t,X(t), m\frac{d}{dt}X(t);\gamma \frac{d}{dt}X(t)+m\frac{d^2}{dt^2}X(t)\right)dt \biggr]:\right.\nonumber\\
\label{1.10}\\
&&\qquad\left.\frac{d}{dt}X(\cdot)\in AC([0,1];\mathbb{R}^d), {\rm a.s.}, 
P^{(X(0), X(1))}\in \Pi(P_0,P_1)\right\}.\nonumber
\end{eqnarray}
If $\sigma(\cdot,\cdot)=0$ and $L=|u|^2$, then for $m>0$,
\begin{equation}
V^m(P_0,P_1)\le \gamma^2 T(P_0,P_1),
\end{equation}
where the equality holds if we formally substitute $m=0$ in (\ref{1.10}) (see (\ref{0.1})--(\ref{1.2.0}) and also Remark \ref{rk1.1.0}).
Indeed,
\begin{eqnarray*}
&&\int_0^1 \left|\gamma \frac{d}{dt}X(t)+m\frac{d^2}{dt^2}X(t)\right|^2dt\\
&=&\int_0^1 \left\{\gamma^2\left| \frac{d}{dt}X(t)\right|^2+m^2\left| \frac{d^2}{dt^2}X(t)\right|^2\right\}dt+
\gamma  m\left(\left| \frac{d}{dt}X(1)\right|^2-\left| \frac{d}{dt}X(0)\right|^2\right).
\end{eqnarray*}
\end{remark}




We consider the case where $m=0$.
Substitute $m=0$ in %(\ref{1.1}) and 
(\ref{1.6}).
%replace $Y(t)dt$ by $mdX(t)$ and then substitute $Y(t)dt=0dX(t)=0$ (see (\ref{1.1})).
Then we formally obtain the following SDE:
\begin{eqnarray}
%Y(t)dt&=&0dX(t)=0,\nonumber\\
\gamma dX(t)&=& u_X(t)dt+\sigma_0(t,X(t))dW(t),\quad 0< t<1,\label{1.4}%\\
%P^{X(t)}&=&P_t, \quad t=0,1.
\end{eqnarray}
where
$$\sigma_0(t,x):=\sigma (t,(x,0)), \quad (t,x,0)\in [0,1]\times \mathbb{R}^d\times \mathbb{R}^d.$$
We denote by $\mathcal{A}$ the set of $\mathbb{R}^{d}$--valued continuous semimartingales $\{X(t)\}_{0\le t\le 1}$ defined on a complete filtered probability space such that (\ref{1.4}) holds weakly.
Let 
\begin{eqnarray*}
L_0(t,x;u):&=&L(t,x,0;u), \quad (t,x,0,u)\in [0,1]\times \mathbb{R}^d \times\mathbb{R}^d \times \mathbb{R}^d,\\
\mathcal{A}(P,Q):&=&\{\{ X(t)\}_{0\le t\le 1}\in \mathcal{A}:P^{(X(0), X(1))}\in \Pi(P,Q)\},\quad P,Q\in \mathcal{P}(\mathbb{R}^d).
\end{eqnarray*}
The following is an SOT for the Langevin dynamics with mass $m=0$,
and is a well--known class of SOTs (see e.g., \cite{M08}--\cite{MT06} and the references therein).

\begin{definition}%[Stochastic optimal transport\label{def1}]

\noindent
For $P_0,P_1\in \mathcal{P}(\mathbb{R}^d )$, let
\begin{equation}\label{1.12}
V^0(P_0,P_1)
:=\inf \left\{E\biggl[\int_0^1 L_0(t,X(t);u_X (t))dt \biggr]:X\in \mathcal{A}(P_0,P_1)\right\}.
\end{equation}
If the set over which the infimum is taken is empty, we set the infimum to be infinite.
\end{definition}

\begin{remark}
If $\sigma_0=0$ and $L=|u|^2$, then $V^0(P_0,P_1)=\gamma^2T(P_0,P_1)$
(see also Remarks \ref{rk1.1.0} and \ref{rk1.3}).

\end{remark}


%Substitute $m=0$ in Remark \ref{rk1.3} and 
%compare $V^0(P_0,P_1)$ with $T(P_0,P_1)$ (see (\ref{0.1}) and also 
%\cite{L2,M04,M09,M2021,1-MTsiam}).

For $m>0$ and a closed set $B\subset \mathcal{P}(\mathbb{R}^d)$, 
we show the existence of  minimizers of $V^m(B,P_0;P_1)$ when it is finite.
Let $\delta_0(dy)$ denote the delta measure on $\{0\}\subset \mathbb{R}^d$.
%and $m_n>0, n\ge 1$ and $m_n\downarrow 0$ as $n\to\infty$.
We also show that if $D^m\subset \mathcal{P}(\mathbb{R}^d)$ ``converges'' to $\delta_0$
as $m\to 0$, then
%there exist closed sets $\{B_n\}_{n\ge 1}\subset\mathcal{P}(\mathbb{R}^d)$ such that $\cap_{n\ge 1}B_n=
%\{\delta_0\}$ and such that 
$V^{m}(D^m,P_0;P_1)$ and its minimizer converge to $V^0(P_0,P_1)$ and its minimizer as $m\to 0$ 
%This implies that $V^0(P_0,P_1)$ is finite if and only if $V^m(P_0,P_1)$ is bounded for sufficiently small $m>0$
(see section \ref{sec:2} for an exact meaning of ``converge''). 

We will study the duality formula for $V^m(B,P_0;P_1)$ for $m>0$ somewhere else
(see Propositions \ref{pp2.2}--\ref{pp2.3} and Example \ref{ex2.1} in section \ref{sec:2} for some discussion).

In section \ref{sec:2}, we state our results.
In section \ref{sec:3}, we give technical lemmas.
In section \ref{sec:4}, we prove our results.














%\newpage

\section{Main results}\label{sec:2}


In this section, we state our results.
We discuss the zero--mass limit of the SOTs when $L(t,z;u)$ is not necessarily of polynomial growth in $u$ and when $L(t,z;u)$ is of polynomial growth in $u$.
We also discuss the duality formula for the SOT.

We state all assumptions before we state our results.
Let 
\begin{equation}
C_{r,R}:=\inf\left \{\frac{L(t,z;u)}{|u|^r}: t\in [0,1], z\in \mathbb{R}^{2d},u \in \mathbb{R}^d, |u|\ge R\right\},\quad r\ge 1, R>0,
\end{equation}
\begin{equation}
R_1(t,z;u)=R_1(L)(t,z;u):=L(t,z;u)-L(t,z;0),\quad (t,z,u)\in [0,1]\times \mathbb{R}^{2d}\times\mathbb{R}^d.
\end{equation}
For $\varepsilon_1\ge 0,\varepsilon_2\in (0,\infty]$, let
$$\Delta L(\varepsilon_1,\varepsilon_2):=\sup \frac{L(t_1,z_1;u)-L(t_2,z_2;u)}{1+L(t_2,z_2;u)},$$
where the supremum is taken over all $u\in \mathbb{R}^d $ and 
all $(t_1,z_1)$ and  $(t_2,z_2)\in [0,1]\times \mathbb{R}^{2d}$
for which $|t_1-t_2|\le\varepsilon_1$, $|z_1-z_2|< \varepsilon_2$.


We describe our assumptions.

\noindent
(A1)
(i) $\sigma:[0,1]\times\mathbb{R}^{2d}\rightarrow M(d,\mathbb{R})$ is bounded and continuous.

\noindent
(ii) 
\begin{equation}
\lim_{R\to\infty}C_{1,R}=\infty.
\end{equation}

\noindent
(iii)
$L:[0,1]\times\mathbb{R}^{2d}\times\mathbb{R}^d\rightarrow[0,\infty)$ is lower semicontinuous.

\noindent
(iv) For $(t,z)\in [0,1]\times\mathbb{R}^{2d}$, $L(t,z;\cdot)$ is convex.

\noindent
(A2) (i) $\sigma$ is a $d\times d$--identity matrix.

\noindent 
(ii) $L:[0,1]\times\mathbb{R}^{2d}\times\mathbb{R}^d
\rightarrow[0,\infty)$ is continuous.

\noindent
(iii) 
\begin{equation}\label{2.2}
R_1(t,z;ru)\le r^2R_1(t,z;u),\quad  t\in [0,1], z\in \mathbb{R}^{2d},u\in \mathbb{R}^d, 0<r<1.
\end{equation}
(iv) There exists $\varepsilon_0>0$ such that $\Delta L(\varepsilon_0,\infty)$ is finite.

\noindent
(A3)
(i) $\Delta L(\varepsilon_1,\varepsilon_2)\to
0$, as $\varepsilon_1,$ $\varepsilon_2\to 0$.

\noindent 
(ii)  For $(t,z)\in [0,1]\times\mathbb{R}^{2d}$, $L(t,z;\cdot)$ is strictly convex.

\noindent 
(iii) 
\begin{equation}
\lim_{R\to\infty}C_{2,R}>0.
\end{equation}

\noindent
(A4) 
(i) There exists $C>0$ and $r_0\ge 1$ such that the following holds:
\begin{equation}\label{2.5.0}
L(t,z;u+v)\le L(t,z;u)+C|v|(|u|^{r_0-1}+|v|^{r_0-1}), \quad  t\in [0,1], z\in \mathbb{R}^{2d},u,v \in \mathbb{R}^d.
\end{equation}
(ii)
\begin{equation}\label{2.7.0}
\lim_{R\to\infty}C_{r_0,R}>0.
\end{equation}

\begin{remark}\label{rk2.1}
(i) (A1, iv) and (A2, iii) imply that $D_uL (t,z;0)$ exists and is equal to $0$.
%In particular,
%$$R_1(L)(t,z;u)=L(t,z;u)-L(t,z;0)-\langle D_uL (t,z;0), u\rangle.$$
Indeed, from (\ref{2.2}), for any subgradient $p_0\in \partial_u L (t,z;0)$, 
$$\langle p_0,rp_0\rangle \le L (t,z;rp_0)-L (t,z;0)\le r^2 \{L (t,z;p_0)-L (t,z;0)\},\quad 0<r<1$$
(see e.g., \cite{1-V2} for subgradient and subdifferential).\\
(ii) The assumptions of Corollary \ref{thm2.2}
 are (A1, ii) and  (A2).
The following function satisfies (A1, ii)  and (A2, ii--iv), but not  (A1, iv): for $p\in (0,2)$, 
$$L=L(u)=|u|^2-|u|^p+1,\quad u\in  \mathbb{R}^d.$$
Indeed, $L (0)=1$ is a local maximum since $|u|^2<|u|^p$ if $0<|u|<1$.\\
(iii) (A2, ii, iv) implies that for $u\in \mathbb{R}^{d}$, $L (\cdot,\cdot;u)\in C_b ([0,1]\times\mathbb{R}^{2d})$
since 
$$L (t,z;u)\le L (t,0;u)+\triangle L(0,\infty)(1+L (t,0;u)),\quad (t, z)\in [0,1]\times \mathbb{R}^{2d}.$$
(iv) (A4) implies that $u\mapsto L(t,z;u)$ grows up of order $|u|^{r_0}$, as $|u|\to\infty$.
Indeed, substituting $u=0$ in (\ref{2.5.0}), the following holds:
$$
L(t,z;v)\le L(t,z;0)+C|v|^{r_0}, \quad  t\in [0,1], z\in \mathbb{R}^{2d},v \in \mathbb{R}^d.
$$
From (\ref{2.7.0}), for sufficientll large $R>0$,
$$0<C_{r_0,R}|v|^{r_0}\le L(t,z;v),\quad  t\in [0,1], z\in \mathbb{R}^{2d},
|v|\ge R.$$
$L=|u|^{r_0}$ satisfies (A4). Indeed, for $u,v \in \mathbb{R}^d$,
$$
|u+v|^{r_0}\le (|u|+|v|)^{r_0}\le  |u|^{r_0}+r_0|v|(|u|+|v|)^{r_0-1},
$$
since $[0,\infty)\ni x\mapsto (|u|+x)^{r_0}$ is convex if $r_0\ge 1$.
$$(|u|+|v|)^{r_0-1}\le 
\begin{cases}
2^{r_0-2}(|u|^{r_0-1}+|v|^{r_0-1}),&\quad r_0\ge 2,\\
|u|^{r_0-1}+|v|^{r_0-1},&\quad 1\le r_0<2.
\end{cases}
$$
since $[0,\infty)\ni x\mapsto x^{r_0-1}$ is convex if $r_0\ge 2$,
and since $[0,\infty)\ni x\mapsto (x+|v|)^{r_0-1}-x^{r_0-1}$ is nonincreasing if $1\le r_0<2$.\\
(v) Let $U\in UC_b ([0,1]\times \mathbb{R}^{2d};[0,\infty))$.
Let also $\{a_n,p_n\}_{n\ge 1}$ 
such that $a_n>0$ for at least one $n\ge 1$, that $a_n\ge 0, 2\le p_n<p_{n+1},  n\ge 1$, and such that 
 $$L_1(u):=\sum_{n=1}^\infty a_n |u|^{p_n}<\infty, \quad u\in\mathbb{R}^d.$$
The following is an example of $L$ that satisfies (A1, ii), (A2, ii--iv), and (A3): 
$$L (t,z;u)=L_1(u)+U(t,z),\quad (t, z,u)\in  [0,1]\times  \mathbb{R}^{2d}\times\mathbb{R}^d.$$
%(v) If (A1, iv) and (A2, iii) holds and if  $L(t,z;u)$ has bounded second derivatives, then (A4) holds.
If there exists $n_0$ such that $a_n=0, n\ge n_0$, then
$L$ given above also satisfies  (A4).
%: for $p\ge 2$, $$L (t,z;u)=|u|^p+U(t,z),\quad (t, z,u)\in  [0,1]\times  \mathbb{R}^{2d}\times\mathbb{R}^d.$$
%(v) If (A1, iv) and (A2, iii) holds and if  $L(t,z;u)$ has bounded second derivatives, then (A4) holds.
\end{remark}

\subsection{Zero--mass limit: general cost function}

In this section, we do not assume that $u\mapsto L(t,z;u)$ is of polynomial growth and discuss the zero--mass limit of SOTs.

Modifying the idea in \cite{M08,M21} (see also \cite{M2021}), the following holds.

\begin{theorem}\label{pp1.1}
Suppose that (A1, i, ii) holds and that $m>0$.
Then for any closed set $B\subset \mathcal{P}(\mathbb{R}^d)$, 
any  $P_0,P_1\in\mathcal{P}(\mathbb{R}^d)$, and for any $\{Z_n=(X_n, Y_n)\}_{n\ge 1}\subset \mathcal{A}^m(B,P_0;P_1)$ such that 
\begin{equation}\label{2.4}
\sup_{n\ge 1}E\biggl[\int_0^1 L(t,Z_n(t);u_{X_n} (t))dt \biggr]<\infty,
\end{equation}
$$\{Q_n(dt\hbox{ }dz\hbox{ }du):=dtP^{(Z_n(t),u_{X_n} (t))}(dz\hbox{ }du)\}_{n\ge 1}$$ 
and $\{Z_n\}_{n\ge 1}$ are tight.
For any weak limit point $Q_\infty(dt\hbox{ }dz\hbox{ }du)$ of $\{Q_n(dt\hbox{ }dz\hbox{ }du)\}_{n\ge 1}$,
there exists $Z=(X,Y)\in \mathcal{A}^m(B,P_0;P_1)$ such that
\begin{eqnarray}
u_X(t)&=&E^{Q_\infty}[u|t,Z(t)],\quad dtdP{\rm -a.e.},\label{2.5}\\
Q_\infty(dt\hbox{ }dz\times \mathbb{R}^d)&=&dtP^{Z(t)}(dz),\label{2.6}
\end{eqnarray}
where $E^{Q_\infty}[u|t,z]$ denotes a conditional expectation of $u$ given $(t,z)$ under $Q_\infty$.\\
Suppose, in addition,  that (A1, iii, iv) holds. Then
there exist a weak limit point $Q_\infty$ of $\{Q_n\}_{n\ge 1}$ and $Z=(X,Y)\in \mathcal{A}^m(B,P_0;P_1)$ such that (\ref{2.5})--(\ref{2.6}) hold and such that 
\begin{equation}
\liminf_{n\to\infty}E\biggl[\int_0^1 L(t,Z_n(t);u_{X_n} (t))dt \biggr]
\ge E\biggl[\int_0^1 L(t,Z(t);u_{X} (t))dt \biggr].\label{1.9.1}
\end{equation}
In particular, $V^m(B,P_0;P_1)$ has a minimizer  $Z\in \mathcal{A}^m(B,P_0;P_1)$ such that (\ref{2.5})--(\ref{2.6}) hold 
for some $Q_\infty\in \mathcal{P}([0,1]\times \mathbb{R}^{2d}\times \mathbb{R}^{d})$, provided it is finite.
\end{theorem}

For $m>0$ and $f:%=(f_i)_{i=1}^d$, $f_i:
[0,1]\rightarrow (\mathbb{R}\cup\{\infty\})^d$, let % 1\le i\le d$, let
\begin{eqnarray}\label{2.9}
\Psi^{m} (f)(t):&=&
\int_0^t \frac{\gamma}{m}\exp\left(-\frac{\gamma (t-s)}{m} \right)f(s)ds, %\quad 0\le t\le T,
\end{eqnarray}
for $t\in [0,1]$ such that the r. h. s. of (\ref{2.9}) is well defined.

%and $\Psi^{m} (f):=\Psi^{m} (1,f)$ for $f\in C([0,1];\mathbb{R}^d)$.
%Then $\Psi^{m} $ is Lipschitz continuous on $C([0,1];\mathbb{R}^d)$.
For $m>0$ and $Z=(X,Y)\in \mathcal{A}^m$, let
\begin{eqnarray}
U_X(t):&=&\int_0^t u_X(s)ds,\quad 0\le t\le1,\label{2.7}\\
M_{X}(t):&=&\int_0^t\sigma(s,Z(s))dW(s),\quad 0\le t\le1\label{2.8}
\end{eqnarray}
 (see (\ref{1.2}) for notation and also Remark \ref{rk1.1}).

For $X\in \mathcal{A}$, we also use notations which are similar to (\ref{2.7})--(\ref{2.8}) when it is not confusing (see (\ref{1.4}) for notation).
In particular, $M_{X}$ also denotes a martingale part of $X\in \mathcal{A}$ for which $\sigma(s,Z(s))$ is replaced by $\sigma_0(s,X(s))$ in (\ref{2.8}).

For $m>0$ and $t\in (0,1]$, $\Psi^{m} $ is Lipschitz continuous on $C([0,t];\mathbb{R}^d)$ (see Lemma \ref{lm3.1}) and the  following holds from (\ref{1.1})--(\ref{1.2}) (see Lemma \ref{lm3.2}): for 
 $Z=(X,Y)\in \mathcal{A}^m$ and $t\in [0,1]$,
\begin{eqnarray}%\label{1.19}
X(t)&=&X(0)+\frac{1}{\gamma}\Psi^{m} (U_X+M_X)(t)
+\frac{1}{\gamma}\left\{1-\exp \left(-\frac{\gamma}{m}t\right)\right\}Y(0), \quad \\
Y(t)&=&U_X(t)+M_X(t)-\Psi^{m} (U_X+M_X)(t)+\exp \left(-\frac{\gamma}{m}t\right)Y(0).
\end{eqnarray}
In particular, $(X(0), Y(0), U_X+M_X)\mapsto Z$  is Lipschitz continuous.

Let $d_{wk}$ denote a metric that induces the topology by weak convergence in $\mathcal{P}(\mathbb{R}^d)$, e.g., Prohorov metric.
The following also holds.

\begin{theorem}\label {thm2.1}
Suppose that (A1, i, ii) holds and that $\{m_n\}_{n\ge 1}$ is a sequence of positive real numbers 
that converges to $0$ as $n\to\infty$.
Then for any $\{B_n\}_{n\ge 1}\subset \mathcal{P}(\mathbb{R}^d)$ such that 
$$\lim_{n\to\infty}\left(\sup\{d_{wk}
(\delta_0, P): P\in B_n\}\right)=0,$$
%$B_n$ is closed, $B_n\supset B_{n+1}, n\ge 1$ and $\cap_{n\ge 1}B_n=\{\delta_0\}$, 
any $P_0,P_1\in\mathcal{P}(\mathbb{R}^d)$, 
and for any $\{Z_n=(X_n, Y_n)\in \mathcal{A}^{m_n}(B_n,P_0;P_1)\}_{n\ge 1}$ such that 
(\ref{2.4}) holds,
%$$\sup_{n\ge 1}E\left[\int_0^1 L(t, Z_n(t);u_{X_n}(t))dt\right]<\infty,$$ 
$\{Q_n(dt\hbox{ }dz\hbox{ }du):=dtP^{(Z_n(t),u_{X_n} (t))}(dz\hbox{ }du)\}_{n\ge 1}$ and $\{Z_n\}_{n\ge 1}$ are tight.
For any weak limit point $\overline Q_\infty(dt\hbox{ }dz\hbox{ }du)$ of $\{Q_n(dt\hbox{ }dz\hbox{ }du)\}_{n\ge 1}$,
there exists $X\in \mathcal{A}(P_0,P_1)$ such that
\begin{eqnarray}
u_X(t)&=&E^{\overline Q_\infty}[u|t,X(t)],\quad dtdP{\rm -a.e.},\label{2.13} 
\\
\overline Q_\infty(dt\hbox{ }dz\times \mathbb{R}^d)&=&dtP^{X(t)}(dx)\delta_0(dy).\label{2.14}
\end{eqnarray}
Suppose, in addition,  that (A1, iii, iv) holds. Then
\begin{equation}\label{2.15}
\liminf_{n\to \infty }V^{m_n}(B_n,P_0;P_1)\ge V^0(P_0,P_1).
\end{equation}
\end{theorem}



For $m>0$ and $t\in [0,1]$, let
\begin{eqnarray}
K^m(t):&=&\frac{1}{\gamma}\Psi^{m} (1)(t)=\frac{1}{\gamma}
\left(1- \exp\left(-\frac{\gamma t}{m}\right)\right),\label{3.7}\\
f^m(t):&=&\gamma K^m (1-t)=1- \exp\left(-\frac{\gamma (1-t)}{m}\right),\label{2.12}\\
%F^m(t):&=&\int_0^t f^m(s)^2ds\le t,\label{3.25}\\
\varphi^m(t):&=&1-\int_t^1 f^m(s)^2ds.\label{3.27}
\end{eqnarray}

For $P_0,P_1\in\mathcal{P}(\mathbb{R}^d)$ and $X\in \mathcal{A}(P_0,P_1)$, let $u_X$ and  $W$ satisfy (\ref{1.4}).
For $m>0$, let $W^m$ be a Brownian motion such that 
\begin{equation}\label{2.24.00}
W(\varphi^m(t))-W(\varphi^m(0))=\int_0^t f^m(s)dW^m(s),\quad t\in [0,1]
\end{equation}
by the martingale representation theorem (see e.g., \cite{IW14}).
Then, under (A2, i), the following holds from (\ref{2.12}): for $t\in [0,1]$,
\begin{eqnarray}\label{3.23}
X(\varphi^m(t))-X(\varphi^m(0))
&=&\int_{0}^{t}K^m(1-s)(u_X^m(s)f^m(s)ds +d W^m(s)),\qquad%\nonumber
\end{eqnarray}
where 
\begin{equation}\label{2.26.00}
u_X^m(t):=u_X(\varphi^m(t)),\quad 0\le t\le 1.
\end{equation}
We define $Z^m=(X^m, Y^m)\in \mathcal{A}^{m}$ by the following: for $t\in [0,1]$,
\begin{eqnarray}
X^m(t)&:=&X(0)+\int_0^t \frac{1}{m}Y^m(s)ds, \label{2.21}\\
%u_{X^m}(t)&=&u_X(\varphi^m(t))f^m(t),\quad t\in [0,1],%\label{3.35}\\
Y^m(t)&=&\frac{X(\varphi^m(0))-X(0)}{K^m(1)}
+\int_0^t \left\{u_X^m(s)f^m(s)- \frac{\gamma}{m}Y^m(s)\right\}ds+W^m(t).\nonumber\\
\label{2.24}
\end{eqnarray}
The following implies that $Z^m=(X^m, Y^m)$ converges to $(X, 0)$ as $m\to 0$.

\begin{theorem}\label{thm2.3}
Suppose that (A2, i) holds and that $m>0$.
Then  for any $P_0,P_1\in\mathcal{P}(\mathbb{R}^d)$, 
any $X\in\mathcal{A}(P_0,P_1)$, and for $Z^m=(X^m, Y^m)\in \mathcal{A}^{m}$ defined by (\ref{2.21})--(\ref{2.24}), the following holds:
\begin{equation}\label{2.29.00}
X^m(1)=X(1).
\end{equation}
In particular, $Z^m=(X^m, Y^m)\in \mathcal{A}^{m}(P_0,P_1)$.
Besides, %for $t\in [0,1)$,
%\begin{eqnarray}
%X^m(t)&=&\left(1-\frac{K^m(t)}{K^m(1)}\right)X(0)
%+f^m(t)\Psi^m\left(\frac{X(\varphi^m(\cdot))}{f^m(\cdot)^2}\right)(t),\label{3.26}\\
%Y^m(t)&=&-\frac{\exp(-\gamma t/m)}{K^m(1)}X(0)
%+\frac{X(\varphi^m(t))}{K^m(1-t)}
%-\gamma\Psi^m\left(\frac{X(\varphi^m(\cdot))}{f^m(\cdot)^2}\right)(t),\label{2.27}
%\end{eqnarray}
%and 
$(X^m, Y^m)$ converges to $(X, 0)$, as $m\to 0$, locally uniformly on $[0,1)$, a.s..\\
Suppose, in addition, that (A2, ii--iv) holds and that
\begin{equation}
E\left[\int_0^1 L_0(t,X(t);u_X(t))dt\right]<\infty.
\end{equation}
Then 
\begin{equation}\label{2.29}
\sup_{m\in (0, \varepsilon_0\gamma/2]}E\left[\int_0^1 L(t,Z^m(t);u_X^m(t)f^m(t))dt\right]<\infty,
\end{equation}
and 
\begin{equation}\label{2.30}
\limsup_{m\to 0}E\left[\int_0^1 L(t,Z^m(t);u_X^m(t)f^m(t))dt\right]\le E\left[\int_0^1 L_0(t,X(t);u_X(t))dt\right].
\end{equation}
\end{theorem}


Suppose that (A1, ii) holds.
For $m>0$ and $P_0,P_1\in \mathcal{P}(\mathbb{R}^d )$, let
\begin{equation}\label{2.3300}
C(m,P_0,P_1):=\frac{1}{f^m(0)}\inf\left\{R\varphi^m(0)+\frac{V^0(P_0,P_1)+1}{C_{1,R}
}+\{d\varphi^m(0)\}^{1/2}:R>0\right\}
\end{equation}
(see (\ref{2.12})--(\ref{3.27}) for notation).
If $V^0(P_0,P_1)$ is finite, then,  for $R>0$,
\begin{eqnarray}
C(m,P_0,P_1)
&\le& \frac{1}{f^m(0)}\left\{R\varphi^m(0)+\frac{V^0(P_0,P_1)+1}{C_{1,R}}+\{d\varphi^m(0)\}^{1/2}\right\}\qquad\\
&\to &\frac{V^0(P_0,P_1)+1}{C_{1,R}}, \quad m\to 0\nonumber\\
&\to &0, \quad R\to\infty.\nonumber
\end{eqnarray}
Besides, since $m\mapsto f^m(\cdot)$ is nonincreasing and $m\mapsto \varphi^m (\cdot)$ is nondecreasing,
\begin{equation}
C(m,P_0,P_1)\downarrow 0, \quad m\downarrow 0.
\end{equation}

Let 
$$\mathcal{P}_{1,C}(\mathbb{R}^d ):=\left\{P\in \mathcal{P}(\mathbb{R}^d ):\int_{\mathbb{R}^d}|x|P(dx)\le C\right\}, \quad  C\ge 0.$$
Then $\mathcal{P}_{1,C}(\mathbb{R}^d )$ is a closed subset of $\mathcal{P}(\mathbb{R}^d )$.

The following holds from Theorem \ref{thm2.3}.

\begin{corollary}\label{thm2.2}
Suppose that (A1, ii) and (A2) hold. Then
for any $P_0,P_1\in \mathcal{P}(\mathbb{R}^d )$ such that $V^0(P_0,P_1)$ is finite
and any $\{D^m\}_{m\in (0, \varepsilon_0\gamma/2]}$ such that 
\begin{equation}\label{2.36.0}
\mathcal{P}_{1,C(m,P_0,P_1)}(\mathbb{R}^d ) \subset D^m, \quad m\in (0, \varepsilon_0\gamma/2],
\end{equation}
$\{V^m(D^m,P_0;P_1)\}_{m\in (0, \varepsilon_0\gamma/2]}$ is bounded and the following holds: 
\begin{equation}\label{2.26}
\limsup_{m\to 0}V^m(D^m,P_0;P_1)\le V^0(P_0,P_1).
\end{equation}
\end{corollary}

\begin{remark}\label{rk2.5}
For $m>0$ and $P_0,P_1\in \mathcal{P}(\mathbb{R}^d )$,
if $D_1\subset D_2\subset \mathcal{P}(\mathbb{R}^d )$, then 
\begin{equation}
V^m(D_2,P_0;P_1)\le V^m(D_1,P_0;P_1).
\end{equation}
\end{remark}

The following holds immediately from Theorem \ref {thm2.1} and Corollary \ref{thm2.2}.
We omit the proof.

\begin{corollary}\label{co2.1}
Suppose that (A1, ii, iv) and (A2) hold.
% and that $\{m_n\}_{n\ge 1}$ is a sequence of positive real numbers such that $m_n\to 0$ as $n\to\infty$.
Then for any $P_0,P_1\in\mathcal{P}(\mathbb{R}^d)$ and 
any $\{D^m
\}_{m\in (0, \varepsilon_0\gamma/2]}\subset \mathcal{P}(\mathbb{R}^d)$ such that (\ref{2.36.0}) and 
\begin{equation}\label{2.39}
\sup\{d_{wk}(\delta_0, P): P\in D^m\}\to 0, \quad m\to 0,
\end{equation}
the following holds:
\begin{equation}
\lim_{m\to  0}V^{m}(D^m,P_0;P_1)=V^0(P_0,P_1).
\end{equation}
\end{corollary}

From Theorem \ref{pp1.1} and Corollary \ref{co2.1} stated above, 
Lemmas  \ref{lm3.3}--\ref{lm3.4} in section \ref{sec:3}, and Theorem 2.4 in \cite{2-JS}, p. 528, 
the first part of the following corollary can be proven in the same way as 
Lemma 3.1 in \cite{MT06}.
Indeed, 
%by the following inequality and Fatou's lemma: 
for $Z=(X,Y)\in \mathcal{A}^m$ and $n\ge 2, \delta >0$, if 
$$\sup\left\{|Z(t)-Z(s)|: t,s\in [0,1], |t-s|\le \frac{1}{n}\right\}<\delta,$$
then under (A1, iv), by Jensen's inequality,
\begin{eqnarray}
&&\int_{1/n}^1
L\left(t,Z(t);n\int_{t-1/n}^t u_X(s)ds\right)dt\\
&\le& \int_{1/n}^1dt
\int_{t-1/n}^t nL\left(t,Z(t);u_X(s)\right)ds\nonumber\\
&\le& \left(1+\Delta L\left(\frac{1}{n}, \delta\right)\right)
\int_0^1 L\left(t,Z(t);u_X(t)\right)dt+\Delta L\left(\frac{1}{n}, \delta\right).\nonumber
\end{eqnarray}
The second part can also be  proven in the same way as  Proposition 2.2, (ii) in \cite{MT06}
(see also \cite{M21}, Proposition 1, (iii)).
Indeed, (A3, ii, iii) implies the uniqueness of the minimizer of $V^0(P_0,P_1)$.
We omit the proof.


\begin{corollary}\label {co2.3}
Suppose that (A1, ii, iv), (A2), and (A3, i) hold.
% and that $\{m_n\}_{n\ge 1}$ is a sequence of positive real numbers 
%such that $(0, \varepsilon_0\gamma/2]\ni m_n\to 0$ as $n\to\infty$.
Then for any $P_0,P_1\in \mathcal{P}(\mathbb{R}^d )$ such that $V^0(P_0,P_1)$ is finite
and any $\{D^m\}_{m\in (0, \varepsilon_0\gamma/2]}\subset \mathcal{P}(\mathbb{R}^d)$ such that $D^m, m\in (0, \varepsilon_0\gamma/2]$ are closed and that (\ref{2.36.0}) and (\ref{2.39}) hold,
any weak limit point, as $m\to 0$, of  any minimizers of $V^{m}(D^m
, P_0;P_1)$ can be written as $(X, 0)$, where $X$ is a minimizer of $V^0(P_0,P_1)$.
Suppose, in addition, that (A3, ii, iii) holds. 
Then minimizers of $V^{m}(D^m, P_0;P_1)$ weakly converge to $(X, 0)$ as $m\to 0$,
where $X$ is the unique minimizer of $V^0(P_0,P_1)$.
%$\sup\{d_{wk}
%(\delta_0, P): P\in B_n\}\to 0, n\to\infty$ and that 
%$B_n\supset \mathcal{P}_{1,C(m_n,P_0,P_1)}(\mathbb{R}^d)$, $n\ge 1$,
%and for any minimizers $Z_n=(X_n, Y_n)$ of $V^{m_n}(B_n, P_0;P_1)$,
%any weak limit point, as $n\to \infty$, of $X_n$ is a minimizer of $V^0(P_0,P_1)$.
%Suppose, in addition, that (A3, ii, iii) holds. 
%Then, for any minimizers $Z_n=(X_n, Y_n)$ of $V^{m_n}(B_n, P_0;P_1)$,
%$X_n$ weakly converges to the unique minimizer of $V^0(P_0,P_1)$, as $n\to \infty$.
\end{corollary}



















\subsection{Zero--mass limit: cost function of polynomial growth in $u$}

In this section, we study the zero--mass limit of SOTs in the case where $u\mapsto  L(t,z;u)$ is of  polynomial growth (see Remark \ref{rk2.1}, (iv)).

For $\{m_n\}_{n\ge 1}$ such that $0<m_n\to 0, n\to\infty$ and $P_0,P_1\in\mathcal{P}(\mathbb{R}^d)$, and for $\{(X_n, Y_n)\in \mathcal{A}^{m_n}(P_0,P_1)\}_{n\ge 1}$ such that $\{E[dX_n(t)/dt|_{t=0}]\}_{n\ge 1}$ is bounded, 
$$\{P^{Y_n(0)}\}_{n\ge 1}\subset\mathcal{P}_{1,C(m_n,P_0,P_1)}(\mathbb{R}^d)$$ for sufficiently large $n\ge 1$.
Indeed, from (\ref{2.3300}),
\begin{equation}
\frac{C(m,P_0,P_1)}{\sqrt m}
\ge\frac{\{d\varphi^m(0)\}^{1/2}}{f^m(0)\sqrt m} 
\to \sqrt{\frac{3d}{2\gamma}},\quad m\to 0.
\end{equation}
We show that (\ref{2.26}) holds under additional assumption (A4),
even when $E[dX(t)/dt|_{t=0}]$ is bounded for $(X,Y)$ under consideration.
%Let
%\begin{equation}
%\mathcal{P}_{2}(\mathbb{R}^d):=\left\{P\in \mathcal{P}_{2}(\mathbb{R}^d):\int_{\mathbb{R}^d}|x|^2P(dx)<\infty\right\}.
%\end{equation}


Under (A4),
for $m>0$, $P_0,P_1\in\mathcal{P}(\mathbb{R}^d)$, and $X\in \mathcal{A}(P_0,P_1)$, 
taking a different probability space, e.g., a product probability space,  if necessary, take 
$\mathcal{Y}^m(0)$, defined on the same probability space as $\{X(t)\}_{0\le t\le 1}$,
that is independent of $X$ and 
such that as $m\to 0$,
\begin{equation}\label{2.42.0}
\mathcal{Y}^m(0)\to 0, \quad {\rm a.s.},\quad 
E[|\mathcal{Y}^m(0)|^{r_0}]\to 0.
\end{equation}


We define $\mathcal{Z}^m=(\mathcal{X}^m, \mathcal{Y}^m)\in \mathcal{A}^{m}$ by the following: for $t\in [0,1]$,
\begin{eqnarray}
\mathcal{X}^m(t)&:=&X(0)+\int_0^t \frac{1}{m}\mathcal{Y}^m(s)ds, \label{2.42}
\\
%u_{X^m}(t)&=&u_X(\varphi^m(t))f^m(t),\quad t\in [0,1],%\label{3.35}
\mathcal{Y}^m(t)&=&\mathcal{Y}^m(0)+
\int_0^t \left\{\left(u_{X}^m(s)+\beta^m(0)\right)f^m(s)- \frac{\gamma}{m}\mathcal{Y}^m(s)\right\}ds
+W^m(t),\nonumber\\
\label{2.43}
\end{eqnarray}
where
\begin{equation}\label{2.44}
\beta^m(0):=\frac{\gamma(X(\varphi^m(0))-X(0)-K^m(1)\mathcal{Y}^m(0))}{1-\varphi ^m(0)}
\end{equation}
(see (\ref{3.7})--(\ref{2.24.00}) and (\ref{2.26.00}) for notation).

The following implies that $\mathcal{Z}^m=(\mathcal{X}^m, \mathcal{Y}^m)$ 
converges to $(X, 0)$ as $m\to 0$.

\begin{theorem}\label{thm2.6}
Suppose that (A2, i) holds and that $m>0$.
Then  for any $P_0,P_1\in\mathcal{P}(\mathbb{R}^d)$, any $X\in\mathcal{A}(P_0,P_1)$, and 
any $\mathcal{Y}^m(0)$ that is independent of $X$ and such that (\ref{2.42.0}) holds, and for $\mathcal{Z}^m=(\mathcal{X}^m, \mathcal{Y}^m)\in \mathcal{A}^{m}$ defined by (\ref{2.42})--(\ref{2.43}), the following holds:
\begin{equation}
\mathcal{X}^m(1)=X(1).
\end{equation}
In particular, $\mathcal{Z}^m=(\mathcal{X}^m, \mathcal{Y}^m)\in \mathcal{A}^{m}(P_0,P_1)$.
Besides,
$(\mathcal{X}^m,\mathcal{Y}^m)$ converges to $(X,0)$, as $m\to 0$, locally uniformly on $[0,1)$, a.s..\\
Suppose, in addition, that (A2, ii--iv), and (A4) hold and that
\begin{equation}
E\left[\int_0^1 L_0(t,X(t);u_X(t))dt\right]<\infty.
\end{equation}
Then 
\begin{equation}\label{2.48}
\sup_{m\in (0, \varepsilon_0\gamma/2]}E\left[\int_0^1 L(t,Z^m(t);(u_X^m(t)+\beta^m(0))f^m(t))dt\right]<\infty,
\end{equation}
and 
\begin{equation}\label{2.49}
\limsup_{m\to 0}E\left[\int_0^1 L(t,Z^m(t);(u_X^m(t)+\beta^m(0))f^m(t))dt\right]\le E\left[\int_0^1 L_0(t,X(t);u_X(t))dt\right].
\end{equation}
\end{theorem}

The following easily holds from Theorem \ref{thm2.6}.
Indeed, if $\{Y_n\}_{n\ge 1}$ is $L^{r_0}$--convergent, then taking a different probability space if necessary,
one can assume that $\{Y_n\}_{n\ge 1}$ is convergent a.s. by Skorokhod's theorem.
We omit the proof.

\begin{corollary}\label{co2.4}
Suppose that (A2) and (A4)  hold. Then
for any $P_0,P_1\in \mathcal{P}(\mathbb{R}^d )$ such that $V^0(P_0,P_1)$ is finite
and any $\{D^m \}_{m\in (0, \varepsilon_0\gamma/2]}\subset \mathcal{P}(\mathbb{R}^d)$ such that 
\begin{equation}\label{2.52}
\inf\{E[|Y|^{r_0}]:P^Y\in D^m\}\to 0, \quad m\to 0,
\end{equation}
$\{V^m(D^m,P_0;P_1)\}_{m\in (0, \varepsilon_0\gamma/2]}$ is bounded and the following holds: 
\begin{equation}\label{2.50}
\limsup_{m\to 0}V^m(D^m
,P_0;P_1)\le V^0(P_0,P_1).
\end{equation}
\end{corollary}

The following holds from Theorem \ref {thm2.1} and Corollary \ref{co2.4}.
We omit the proof.

\begin{corollary}\label {co2.5}
Suppose that (A1, iv), (A2), and (A4) with $r_0>1$ hold.
Then for any $P_0,P_1\in\mathcal{P}(\mathbb{R}^d)$ and 
any $\{D^m\}_{m\in (0, \varepsilon_0\gamma/2]}\subset \mathcal{P}(\mathbb{R}^d)$ such that 
(\ref{2.39}) and (\ref{2.52}) hold,
\begin{equation}
\lim_{n\to \infty}V^{m}(D^m,P_0;P_1)=V^0(P_0,P_1).
\end{equation}
\end{corollary}

From Theorem \ref{pp1.1} and Corollary \ref{co2.5} stated above, 
Lemmas  \ref{lm3.3}--\ref{lm3.4} in section \ref{sec:3}, and Theorem 2.4 in \cite{2-JS}, p. 528, 
the following holds in the same way as Corollary \ref{co2.3}.
We omit the proof.

\begin{corollary}\label {co2.6}
Suppose that (A1, iv), (A2), (A3, i), and (A4) with $r_0>1$ hold.
Then for any $P_0,P_1\in \mathcal{P}(\mathbb{R}^d )$ such that $V^0(P_0,P_1)$ is finite and 
any $\{D^m\}_{m\in (0, \varepsilon_0\gamma/2]}\subset \mathcal{P}(\mathbb{R}^d)$ 
such that $D^m, m\in (0, \varepsilon_0\gamma/2]$ are closed and
 that (\ref{2.39}) and (\ref{2.52}) hold,
any weak limit point, as $m\to 0$, of
any minimizers of $V^m(D^m,P_0;P_1)$ can be written as $(X, 0)$, where $X$ is a minimizer of
$V^0(P_0,P_1)$. 
Suppose, in addition, that (A3, ii, iii) holds. Then minimizers of $V^m(D^m,P_0;P_1)$
weakly converge to $(X, 0)$ as $m\to 0$, where $X$ is the unique minimizer of $V^0(P_0,P_1)$.
\end{corollary}


%For $m>0, \mu,\nu\in \mathcal{P}(\mathbb{R}^d \times \mathbb{R}^d )$, let
%\begin{equation}
%{\bf V}^m (\mu,\nu):=\inf \left\{E\biggl[\int_0^1 L(t,Z(t);u_X (t))dt \biggr]:Z\in \mathcal{A}^m,
%P^{Z(0)}=\mu, P^{Z(1)}=\nu\right\}.
%\end{equation}
%Then for $P_0,P_1\in\mathcal{P}(\mathbb{R}^d)$ and $B\subset \mathcal{P}(\mathbb{R}^d)$,
%$$V^m(B,P_0;P_1)=\inf \left\{{\bf V}^m (\mu,\nu):\mu_1=P_0, \mu_2\in B, \nu_1=P_1\right\}.
%$$
%The duality theorem for ${\bf V}^m$ can be proven in the same way as \cite{M21}.














\subsection{Duality formula}
In this section, we discuss the duality formula for $V^m(B, P_0;P_1)$.


By the convex duality, the following holds  in the same way as \cite{M21} (see also \cite{M2021,MT06} and the reference therein).
We give the proof for completeness.

\begin{proposition}\label{pp2.2}
Suppose that (A1) holds and that $m>0$.
Then  for any closed convex subset $B\subset \mathcal{P}(\mathbb{R}^d)$
and any $P_0,P_1\in \mathcal{P}(\mathbb{R}^d )$, 
\begin{equation}
V^m(B, P_0;P_1)
=\sup\left \{\int_{\mathbb{R}^d}f(x)P_1(dx)-V^m(B, P_0;\cdot)^*(f):f\in C_b (\mathbb{R}^d)\right\},
\end{equation}
where 
\begin{eqnarray*}
V^m(B, P_0;\cdot)^*(f)
&:=&\sup\left\{\int_{\mathbb{R}^d}f(x)P(dx)
-V^m(B,P_0;P):P\in  \mathcal{P}(\mathbb{R}^d)\right\}\\
&=&\sup \left\{E\biggl[f(X(1))-\int_0^1 L(t,Z(t);u_X (t))dt \biggr]:\right.\\
&&\qquad \left.Z\in \mathcal{A}^m, 
P^{X(0)}=P_0, P^{Y(0)}\in B\right\}.
\end{eqnarray*}
\end{proposition}

The following proposition explains why we have to consider a family $\{V^m(B, P_0;P_1), B\subset \mathcal{P}(\mathbb{R}^d)\}$ of SOTs when we consider the zero--mass limit.

\begin{proposition}\label{pp2.3}
Suppose that (A2,i) holds and that $m>0$ and $L=L(t;u)$, $(t,u)\in [0,1]\times \mathbb{R}^d$.
Then for any $P_0\in  \mathcal{P}(\mathbb{R}^d)$ and any $f\in  C_b (\mathbb{R}^d)$,
\begin{eqnarray}
&&V^m(\mathcal{P}(\mathbb{R}^d ), P_0;\cdot)^*(f)\label{2.34.1}\\
&=&\sup\left\{E\biggl[f(\eta (1))-\int_0^1 L(t;u_\eta (t))dt \biggr]:
%\label{2.34.1}\\&&\qquad 
d\eta (t) =K^m(1-t)(u_\eta(t)dt+dW(t))\right\}\notag
\end{eqnarray}
(see (\ref{3.7}) for notation), which does not depend on $P_0$,
where $u_\eta(\cdot)$ is a progressively measurable stochastic process defined on the same filtered probability space as $W(\cdot)$.
In particular, $V^m(P_0,P_1)$ does not necessarily converge to $V^0(P_0,P_1)$ as $m\to 0$.
\end{proposition}


Let 
\begin{eqnarray}
H(t,z;p)&:=&\sup\{\langle p,u\rangle -L(t,z;u):u\in\mathbb{R}^d\},\quad (t,z,p)\in [0,1]\times \mathbb{R}^{2d}\times\mathbb{R}^{d},\nonumber\\
\\
a(t,z)&:=&\sigma(t,z)\sigma(t,z)^*,\label{2.56}
\end{eqnarray}
where $\sigma(t,z)^*$ denotes the transpose of $\sigma(t,z)$.

For $f\in C_b(\mathbb{R}^d)$,
let $\psi^m(t,z)=\psi^m(t,z;f)$ be a classical solution of the following,
provided it exists (see e.g., \cite{M2021}): for $(t,z=(x,y))\in (0,1)\times \mathbb{R}^d\times\mathbb{R}^d$,
\begin{eqnarray}
0&=&\frac{\partial\psi^m(t,z)}{\partial t}+\frac{1}{2}Trace(a(t,z)D^2_y\psi^m(t,z))\label{2.33.0}\\
&&\qquad 
+\left\langle \frac{1}{m} D_x\psi^m(t,z)-\frac{\gamma}{m}D_y\psi^m(t,z), y\right\rangle+H(t,z;D_y\psi^m(t,z)),\nonumber\\
\quad \psi^m(1,z)&=&f(x).\label{2.34.0}
\end{eqnarray}

In the following example, we show that the following holds:
\begin{equation}\label{2.33.-1}
V^m(\mathcal{P}(\mathbb{R}^d ),P_0;\cdot)^*(f)=\sup\{\psi^m(0,x,y;f):y\in\mathbb{R}^d \},
\end{equation}
which does not depend on either $P_0\in \mathcal{P}(\mathbb{R}^d)$ or $x\in  \mathbb{R}^d$.

\begin{example}\label{ex2.1}
Suppose that (A2, i) holds and that $L=|u|^2/2, u\in\mathbb{R}^d$. 
For $f\in C_b^2 (\mathbb{R}^d)$ and $x\in \mathbb{R}^d$, 
let 
\begin{equation}
\phi(t,x):=
\begin{cases}
\displaystyle \log \int_{\mathbb{R}^d}\exp (f(y)) \frac{\gamma^d}{\sqrt{2\pi (1-t)}^d}
\exp \left(-\frac{\gamma^{2}|x-y|^2}{2(1-t)}\right)dy,&t\in [0,1),\\
f(x),&t=1.
\end{cases}\label{2.33}
\end{equation}
%When necessary, we use $\phi(t,x;f)$ instead of $\phi(t,x)$.
Then $\phi\in C_b^{1,2}([0,1]\times \mathbb{R}^d)$, and 
$\phi(t,x)=\phi(t,x;f)$ is a classical solution of the following:
\begin{eqnarray}
\frac{\partial\phi(t,x)}{\partial t}+\frac{1}{2\gamma^2}\triangle_x\phi(t,x)+\frac{1}{2\gamma^2}|D_x\phi(t,x)|^2&=&0,
 \quad (t,x)\in (0,1)\times\mathbb{R}^d,\nonumber\\
 \label{2.37}\\
 \phi(1,x)&=&f(x), \quad x\in \mathbb{R}^d.
\end{eqnarray}
Let
\begin{equation}
%M^m(t,z)&:=x+K^m(1-t)y,\quad t\in [0,1],\label{2.34}\\
%A^m (t)&:=1-\int_t^1 K^m(1-s)^2ds,\quad t\in [0,1],\label{2.35}\\
%\int_0^{1-t}\{1-\exp (-(1-t-s))\}^2ds,\quad t\in [0,1],\\
\psi^m(t,z):=\phi(\varphi^m (t),x+K^m(1-t)y),\quad t\in [0,1], z=(x,y)\in \mathbb{R}^d\times \mathbb{R}^d\label{2.36}
\end{equation}
(see (\ref{3.7})--(\ref{3.27}) for notation).
%When necessary, we use $\psi^m(t,z;f)$ instead of $\psi^m(t,z)$.
Then  from (\ref{2.33})--(\ref{2.36}), $\psi^m(\cdot,\cdot)=\psi^m(\cdot,\cdot;f)\in C^{1,2} ([0,1]\times \mathbb{R}^{2d})$, satisfies (\ref{2.33.0})--(\ref{2.34.0}) with $a=$ an identity matrix, $H=H(p)=|p|^2/2$, and 
\begin{equation}\label{2.41}
\sup\{\psi^m(0,x,y;f):y\in\mathbb{R}^d \} 
=\sup\{\phi(\varphi^m (0),y;f):y\in\mathbb{R}^d \},\quad x\in\mathbb{R}^d.
\end{equation}
%since
%$$\psi^m(0,z) =\phi(\varphi^m (0),x+K^m(1)y).$$

We show that (\ref{2.33.-1}) holds. 
%\begin{equation}\label{5.9}
%V^m(P,\cdot)^*(f)=\sup\{\phi(\varphi^m (0),y):y\in\mathbb{R}^d \},
%\end{equation}
For $Z\in \mathcal{A}^m$,  
\begin{eqnarray}\label{5.10}
&&E[f(X(1))]-E\left[\int_0^1 \frac{1}{2}|u_X(t)|^2dt\right]\\
&\le& E[\phi(\varphi^m (0),X(0)+K^m(1)Y(0);f)]\le\sup\{\phi(\varphi^m (0),y;f):y\in\mathbb{R}^d \},\nonumber
\end{eqnarray}
where the equality holds in the first inequality if $D_y\psi^m(t,Z(t))=u_X(t)$.
Indeed, from (\ref{2.33.0}), by the It\^o formula, 
\begin{eqnarray*}
&&d\psi^m(t,Z(t))\\
&=&\left(\frac{\partial\psi^m(t,Z(t))}{\partial t}+\frac{1}{2}\triangle_y\psi^m(t,Z(t))
+\left\langle \frac{1}{m}Y(t),D_x\psi^m(t,Z(t))\right\rangle\right.\\
&&\qquad \left.+\left\langle u_X(t)-\frac{\gamma}{m}Y(t),D_y\psi^m(t,Z(t))\right\rangle\right)dt\\
%&&\qquad +\left.\left\langle \frac{1}{m}D_x\psi^m(t,Z(t))+u_X(t)-\frac{\gamma}{m}D_y\psi^m(t,Z(t)), Y(t)\right\rangle\right)dt\\
&&\qquad +\langle D_y\psi^m(t,Z(t)), dW(t)\rangle\\
&=&\left(-\frac{1}{2}|D_y\psi^m(t,Z(t))|^2+\langle u_X(t), D_y\psi^m(t,Z(t))\rangle\right)dt
+\langle D_y\psi^m(t,Z(t)), dW(t)\rangle\\
&=&\left(-\frac{1}{2}|D_y\psi^m(t,Z(t))-u_X(t)|^2+\frac{1}{2}|u_X(t)|^2\right)dt
+\langle D_y\psi^m(t,Z(t)), dW(t)\rangle.
\end{eqnarray*}

For $y\in\mathbb{R}^d$, %if we take $$Y(0)=\frac{y-X(0)}{K^m(1)},$$then
 from (\ref{5.10}),
\begin{eqnarray}\label{2.67}
&&V^m(\mathcal{P}(\mathbb{R}^d ), P_0;\cdot)^*(f)\\
&\ge &\sup \left\{E\biggl[f(X(1))-\int_0^1 \frac{1}{2}|u_X(t)|^2dt \biggr]:
%\right.\\&&\qquad \left.
Z\in \mathcal{A}^m, 
P^{X(0)}=P_0, Y(0)=\frac{y-X(0)}{K^m(1)}\right\}\nonumber\\
&=&\phi(\varphi^m (0),y;f).\nonumber
\end{eqnarray}
Indeed, the following has a unique strong solution:
\begin{eqnarray*}
dX(t)&=&\frac{1}{m}Y(t)dt,\\%\label{1.1}\\
dY(t)&=&\left\{D_y\psi^m(t,Z(t))-\frac{\gamma }{m}Y(t)\right\}dt+dW(t),\quad 0< t<1,%\label{1.2}%\\
%P^{X(t)}&=&P_t, \quad t=0,1.
\end{eqnarray*}
since
$$D^2_y\psi^m(t,z)=K^m(1-t)^2D^2\phi(\varphi^m (t),x+K^m(1-t)y),\quad 
t\in [0,1], z=(x,y)\in \mathbb{R}^d\times \mathbb{R}^d
$$
is bounded.

Since $y$ is arbitrary in (\ref{2.67}), (\ref{2.41})--(\ref{2.67}) imply (\ref{2.33.-1}).
\end{example}


We will discuss a meaningful characterization of  $V^m(B,P_0;\cdot)^*(f)$ somewhere else.





















%\newpage
\section{Lemmas}\label{sec:3}


In this section, we give technical lemmas.

The following lemma will be used in the proofs of Lemmas \ref{lm3.3}, \ref{lm3.4}, and Theorem \ref{thm2.3}.
We give the proof for completeness  (see (\ref{2.9})  for notation).

\begin{lemma}\label{lm3.1}
Let $m>0$.
(i) For $t\in [0,1]$ and $f\in C([0,t];\mathbb{R}^d)$,
\begin{equation}\label{3.1}
\Psi^{m} (f)(t)\le \|f\|_{\infty,t}\left(1- \exp\left (-\frac{\gamma t}{m} \right) \right),
\end{equation}
where $\|f\|_{\infty,t}:=\sup\{|f(s)|:0\le s\le t\}$.
In particular, for  $f\in C([0,1];\mathbb{R}^d)$,
\begin{equation}\label{3.2}
\|\Psi^{m} (f)\|_{\infty}\le \|f\|_{\infty},
\end{equation}
where $\|f\|_{\infty}:=\|f\|_{\infty,1}$. The following also holds:
\begin{equation}\label{3.3.0}
\lim_{\tilde m\to m}\|\Psi^{\tilde m} (f)-\Psi^{m} (f)\|_{\infty}=0, \quad m>0.
\end{equation}
(ii) For $t\in [0,1], \delta\in (0,1)$,  and $f\in C([0,t];\mathbb{R}^d)$, 
%For $f\in C([0,1];\mathbb{R}^d)$, $\delta>0$, and  $t\in (0,1]$,
\begin{eqnarray}\label{3.3}
&&\left|\Psi^{m} (f)(t)-\left(1-\exp\left (-\frac{\gamma t}{m} \right)\right)f(t)\right|\\
& \le& \left\{2\|f\|_{\infty,t}\exp\left (-\frac{\gamma \delta}{m} \right)
+\sup_{t-t\wedge\delta\le s\le t}|f(t)-f(s)|\right\}\left(1-\exp\left (-\frac{\gamma t}{m} \right)\right),\nonumber
\end{eqnarray}
\begin{equation}\label{3.4.0}
\exp\left (-\frac{\gamma t}{m} \right)|f(t)|\le \|f\|_{\infty,t\wedge\delta}+\exp\left (-\frac{\gamma \delta}{m} \right)\|f\|_{\infty,t},
\end{equation}
where $t\wedge\delta:=\min (t,\delta )$.
In particular, for $f\in C([0,1];\mathbb{R}^d)$ such that $f(0)=0$,
\begin{equation}\label{3.4}
\|\Psi^{m} (f)-f\|_{\infty}\le 3\|f\|_{\infty}\exp\left (-\frac{\gamma \delta}{m} \right)
+2\sup_{t,s\in [0,1], |t-s|\le\delta}|f(t)-f(s)|\to 0,
\end{equation}
as $m\to 0$ and then $\delta \to 0$.
%\begin{equation}\label{3.5}
%\sup_{0<t\le 1}\left|\left(1-\exp\left (-\frac{\gamma t}{m} \right)\right)^{-1}\Psi^{m} (f)(t)-f(t)\right |\to 0, \quad m\to 0.
%\end{equation}

%where $\|f\|_{t, \infty}:=\sup\{|f(s)|:0\le s\le t\}$.
\end{lemma}

\begin{proof}
(\ref{3.1}) is true, since 
%\begin{equation}
%|\Psi^{m} (f)(t)|\le \|f\|_{\infty,t}\int_0^t \frac{\gamma}{m}\exp\left(-\frac{\gamma (t-s)}{m} \right)ds,
%\end{equation}
\begin{equation}\label{3.6}
\int_0^t \frac{\gamma}{m}\exp\left(-\frac{\gamma (t-s)}{m} \right)ds
=\left(1-\exp\left (-\frac{\gamma t}{m} \right)\right).
\end{equation}
 (\ref{3.3.0}) can be proven easily from the following: for $s,t\in [0,1]$ such that $ s\le t$,
\begin{eqnarray*}
&&\left|\frac{\gamma}{m_{n_{k}}}\exp\left(-\frac{\gamma(t-s)}{m_{n_{k}}}\right)-
\frac{\gamma}{m}\exp\left(-\frac{\gamma(t-s)}{m_{n_{k}}}\right)\right|
\le\gamma\left|\frac{1}{m_{n_k}}-\frac{1}{m}\right|, \\
&&\left|\exp\left(-\frac{\gamma(t-s)}{m_{n_{k}}}\right)-
\exp\left(-\frac{\gamma(t-s)}{m}\right)\right|\\
& \le& \exp\left(-\frac{\gamma(t-s)}{\max(m_{n_{k}},m)}\right)
\left(1-\exp\left(-\left|\frac{1}{m}-\frac{1}{m_{n_k}}\right|\gamma(t-s)\right)\right)
\le \left|\frac{1}{m}-\frac{1}{m_{n_k}}\right|\gamma.
\end{eqnarray*}

We prove  (\ref{3.3}).
From (\ref{3.6}),
\begin{eqnarray*}
&&\Psi^{m} (f)(t)-\left(1-\exp\left (-\frac{\gamma t}{m} \right)\right)f(t)\\
&=&\left(\int_0^{t-t\wedge\delta}+ \int_{t-t\wedge\delta}^t\right)\frac{\gamma}{m}
\exp\left(-\frac{\gamma (t-s)}{m} \right)(f(s)-f(t))ds,
\end{eqnarray*}
\begin{eqnarray*}
&&\int_0^{t-t\wedge\delta}\frac{\gamma}{m}\exp\left(-\frac{\gamma (t-s)}{m} \right)|f(s)-f(t)|ds\\
&\le& 2\|f\|_{\infty,t}\exp\left (-\frac{\gamma (t\wedge\delta ) }{m}\right) \left(1-\exp\left (-\frac{\gamma (t-t\wedge\delta)}{m} \right)\right)\\
&&\begin{cases}
=0,&t\wedge\delta=t,\\
\le \displaystyle 2\|f\|_{\infty,t}\exp\left (-\frac{\gamma\delta}{m} \right)\left(1-\exp\left (-\frac{\gamma t}{m} \right)\right),&t\wedge\delta=\delta,
\end{cases}
\end{eqnarray*}
\begin{eqnarray*}
&&\int_{t-t\wedge\delta}^t\frac{\gamma}{m}\exp\left(-\frac{\gamma (t-s)}{m} \right)|f(s)-f(t)|ds\\
& \le& \sup_{t-t\wedge\delta\le s\le t}|f(s)-f(t)|\left(1-\exp\left (-\frac{\gamma (t\wedge\delta )}{m} \right)\right).%\\
%& \le \sup_{0\le s_1, s_2\le 1, |s_1 -s_2|\le t\wedge\delta}|f(s_1)-f(s_2)|\left(1-\exp\left (-\frac{\gamma t}{m} \right)\right).
\end{eqnarray*}
%(\ref{3.3}) implies (\ref{3.5}).
The following implies (\ref{3.4.0}):
\begin{eqnarray*}
\exp\left (-\frac{\gamma t}{m} \right)|f(t)|
\le 
\begin{cases}
|f(t\wedge\delta)|,&t\wedge\delta=t,\\%0\le t\le \delta,\\
\displaystyle \exp\left (-\frac{\gamma \delta}{m} \right)|f(t)|,&t\wedge\delta=\delta.%\le t\le 1.
\end{cases}
\end{eqnarray*}
 \end{proof}

%Let
%\begin{equation}
%K^m(t):=\frac{1}{\gamma}\Psi^{m} (1)(t)=\frac{1}{\gamma}
%\left(1- \exp\left(-\frac{\gamma t}{m}\right)\right),
%\quad  0\le t\le 1.\label{3.7}
%\end{equation}

The following lemma plays a crucial role in the proofs of Lemmas \ref{lm3.3}, \ref{lm3.4}, \ref{lm3.6} and  Theorem \ref{thm2.3} (see (\ref{2.9})--(\ref{2.8}) and (\ref{3.7}) for notation).

\begin{lemma}\label{lm3.2}
For $m>0$, $Z=(X,Y)\in \mathcal{A}^m$, and $t\in[0,1]$,
\begin{eqnarray}
\quad X(t)&=&X(0)+\frac{1}{\gamma}\Psi^{m} (U_X+M_X+Y(0))(t)\label{3.8}\\
%&=&X(0)+K^m(t)Y(0)+\frac{1}{\gamma}\Psi^{m} (U_X+M_X)(t)\label{3.8}\\
&=&X(0)+K^m(t)Y(0)+\int_0^tK^m(t-s)(u_X(s)ds+ \sigma(s,Z(s))dW(s)),\label{3.9}\\
\quad Y(t)&=&\exp \left(-\frac{\gamma}{m}t\right)Y(0)+U_X(t)+M_X(t)-\Psi^{m} (U_X+M_X)(t)\label{3.10}\\
&=&\exp \left(-\frac{\gamma}{m}t\right)Y(0)+\int_{0}^{t}\exp\left(-\frac{\gamma (t-s)}{m}\right)(u_X(s)ds +\sigma(s,Z(s))dW(s)).\nonumber\\
\label{3.11}
\end{eqnarray}
\end{lemma}

\begin{proof}
We  prove (\ref{3.8}).
%\begin{equation}\label{1.16}
%X(t)-X(0)=\int_0^t \frac{1}{m}\exp\left(-\frac{\gamma (t-s)}{m} \right)
%\left(U_X(s)+M_X(s)+mY(0)\right)ds.
%\end{equation}
Integrating (\ref{1.2}) in $t$, from (\ref{1.1}),
%\begin{equation}
%m(Y(t)-Y(0))=U_X(t)-\gamma (X(t)-X(0))+M_X(t),\quad 0\le  s\le t\le 1,\label{1.12}
%\end{equation}
%which implies 
\begin{equation*}
\frac{d}{dt} (X(t)-X(0))=-\frac{\gamma}{m} (X(t)-X(0))+\frac{1}{m}(U_X(t)+M_X(t)+Y(0)),
\end{equation*}
which implies  (\ref{3.8}).

(\ref{3.8}) and the following imply (\ref{3.9}) (see (\ref{3.7}) for notation): by the integration by parts,
\begin{eqnarray}\label{3.11.0}
\frac{1}{\gamma}\Psi^{m} (U_X+M_X)(t)&=&\int_0^t \frac{1}{m}\exp\left(-\frac{\gamma (t-s)}{m} \right)\left(U_X(s)+M_X(s)\right)ds\qquad\\
&=&\int_0^t \left\{-\frac{d}{ds}K^m(t-s)\right\}\left(U_X(s)+M_X(s)\right)ds\nonumber\\
&=&\int_0^t K^m(t-s)\left(u_X(s)ds+ \sigma(s,Z(s))dW(s)\right),\nonumber
\end{eqnarray}
since $K^m(0)=0$ and $U_X(0)=M_X(0)=0$.

Differentiate (\ref{3.8}) and we obtain (\ref{3.10}) from (\ref{1.1}).

From (\ref{3.10}) and (\ref{3.11.0}), we obtain  (\ref{3.11}), since
$$\exp\left(-\frac{\gamma (t-s)}{m}\right)=1-\gamma K^m(t-s),\quad 0\le s\le t\le 1.$$
\end{proof}

The following lemma plays a crucial role in the proofs of Lemma \ref{lm3.4}, Theorem \ref{pp1.1}, and Theorem \ref{thm2.1}.

\begin{lemma}\label{lm3.3}
Suppose that (A1, i, ii) holds and that $\{m_n\}_{n\ge 1}$ is a bounded sequence of positive numbers.
Then for any tight families $\{P_{0,n}\}_{n\ge 1}, \{P_{1,n}\}_{n\ge 1}\subset\mathcal{P}(\mathbb{R}^d)$ and any $\{Z_n=(X_n, Y_n)\in\mathcal{A}^{m_n}(P_{0,n},P_{1,n})\}_{n\ge 1}$ such that (\ref{2.4}) holds,
$\{Q_n (dt\hbox{ }dz\hbox{ }du):=dtP^{(Z_n(t),u_{X_n}(t))}(dz\hbox{ }du)\}_{n\ge 1}$ 
%\subset  \mathcal{P}([0,1]\times \mathbb{R}^{2d}\times \mathbb{R}^{d})$
and $\{(Y_n(0), U_{X_n},M_{X_n})\}_{n\ge 1}$ are tight.
If there exists a positive constant $C$ such that $m_n\ge C, n\ge 1$, 
then $\{Z_n\}_{n\ge 1}$ is tight.
If $m_n\to 0$ as $n\to\infty$, then
$\{Z_n\}_{n\ge 1}$ is tight if and only if 
$$\lim_{n\to\infty}Y_n(0)=0,\quad \hbox{in law}.$$

%\noindent
%(ii)  For $\{X_n\}_{n\ge 1}\subset \mathcal{A}(P_0,P_1)$ such that (\ref{2.4}) with $Y_n(t)$ replaced by $0$ holds,
%\begin{equation*}
%\{Q_n^0 (dt\hbox{ }dx\hbox{ }du):=dtP^{(X_n(t),u_{X_n} (t))}(dx\hbox{ }du)\}_{n\ge 1}\subset  \mathcal{P}([0,1]\times \mathbb{R}^{d}\times \mathbb{R}^{d})
%\end{equation*}
%and $\{(X_n, U_{X_n},M_{X_n})\}_{n\ge 1}$ are tight.

\end{lemma}

\begin{proof}
%Since (ii) can be proven in the same way as (i), we only prove (i) (see also \cite{M08, M21,M2021}).
We first prove that $\{Q_n (dt\hbox{ }dz\hbox{ }du)\}_{n\ge 1}$ is tight.
$[0,1]$ is compact.
For $k\ge 1$ and $R>0$, let 
$$B_{k,R}:=\{z\in \mathbb{R}^{k}: |z|\le R\}.$$ 
Then the following holds uniformly in $n$:
\begin{eqnarray}
Q_n ([0,1]\times B_{2d,R}^c\times \mathbb{R}^{d})&=&\int_0^1 P(|Z_n(t)|> R)dt
\to 0, \quad R\to \infty,\label{1.10.1}\\
Q_n ([0,1]\times \mathbb{R}^{2d}\times B_{d,R}^c)&=&\int_0^1 P(|u_{X_n}(t)|> R)dt\to 0, 
\quad R\to \infty.\qquad\label{1.11.1}
\end{eqnarray}

We first prove (\ref{1.10.1}). For $t\in [0,1]$,
\begin{eqnarray*}
|Z_n(t)|& \le& |X_n(t)|+|Y_n(t)|,\\
|X_n(t)|& \le& |X_n(0)|+ \frac{1}{\gamma}\left(\int_0^1|u_{X_n}(s)|ds
+\sup_{0\le \alpha\le 1}\left|\int_0^\alpha\sigma (s, Z_n(s))dW(s)\right|+|Y_n(0)|\right),\\
|Y_n(t)|& \le& |Y_n(0)|+ 2\left(\int_0^1|u_{X_n}(s)|ds
+\sup_{0\le \alpha\le 1}\left|\int_0^\alpha\sigma (s, Z_n(s))dW(s)\right|\right)
\end{eqnarray*}
from (\ref{3.2}), (\ref{3.8}), and (\ref{3.10}).

$\{|X_n(0)|\}_{n\ge 1}$ is tight since $Z_n=(X_n, Y_n)\in\mathcal{A}^{m_n}(P_{0,n},P_{1,n}), n\ge 1$.

$\{|Y_n(0)|\}_{n\ge 1}$ is also tight from (A1, i, ii) since $Z_n=(X_n, Y_n)\in\mathcal{A}^{m_n}(P_{0,n},P_{1,n}), n\ge 1$.
Indeed, substituting $t=1$ in (\ref{3.8}),
\begin{eqnarray}\label{3.14}
&&\frac{1}{\gamma}\left(1-\exp\left(-\frac{\gamma}{m_n}\right)\right)
|Y_n(0)|\\
&\le& |X_n(1)|+|X_n(0)| +\frac{1}{\gamma}\left(1-\exp\left(-\frac{\gamma}{m_n}\right)\right)\nonumber\\
&&\qquad \times \left(\int_0^1|u_{X_n}(s)|ds
+\sup_{0\le t\le 1}\left|\int_0^t\sigma (s, Z_n(s))dW(s)\right|\right)
\nonumber
\end{eqnarray}
from (\ref{3.7}) and (\ref{3.1}).
The expectation of the third line in (\ref{3.14}) is bounded in $n$, since for  a measurable set $B\subset [0,1]$,
\begin{eqnarray}\label{4.3}
\int_B |u_{X_n}(s)|ds& \le& 
R|B|+\frac{1}{C_{1,R}}\int_0^1 L(t,Z_n(t);u_{X_n} (t))dt,\quad R>0,
\end{eqnarray}
where $|B|$ denote the Lebesgue measure  of $B$, and since
\begin{equation}\label{3.18}
E\left[\sup_{0\le t\le 1}\left|\int_0^t\sigma (s, Z_n(s))dW(s)\right|^2\right]
\le 4E\left[\int_0^1Trace (a (s, Z_n(s)))ds\right]
\end{equation}
by Doob's inequality (see (\ref{2.56}) for notation) .
%Here 
%$$a(s,z):=\sigma (s, z)\sigma  (s, z)^*,$$
%where $\sigma  (s, z)^*$ denotes the transpose of $\sigma  (s, z)$.

The tightness of $\{|X_n(0)|\}_{n\ge 1}$ and $\{|Y_n(0)|\}_{n\ge 1}$ and  (\ref{4.3})--(\ref{3.18}) 
also imply the tightness of $\{\sup_{0\le s\le 1} |X_n(s)|\}_{n\ge 1}$ and
$\{\sup_{0\le s\le 1} |Y_n(s)|\}_{n\ge 1}$, and  (\ref{1.10.1}) holds.

From (\ref{4.3}), (\ref{1.11.1}) can be proven by Chebychev's inequality:
$$\int_0^1 P(|u_{X_n}(t)|> R)dt
\le \int_0^1\frac{1}{R}E[|u_{X_n}(t)|]dt.
$$

From (\ref{4.3})--(\ref{3.18}), and (A1, i), 
$\{( X_n(0),Y_n(0), U_{X_n},M_{X_n}),0\le t\le 1\}_{n\ge 1}$ is also tight % as a sequence of random variables in in $C([0,1];\mathbb{R}^{2d})\times \mathbb{R}^{2d}$ 
 (see \cite{2-JS}, p. 356, Theorem 4.5 and p. 363, Theorem 5.10).
 $\{m_n\}_{n\ge 1}$ is  bounded by assumtion.
Take a convergent subsequence $\{m_{n_k}\}_{k\ge 1}$ and a weakly convergent subsequence 
$\{(X_{n_k}(0), Y_{n_k}(0), U_{X_{n_k}},M_{X_{n_k}})\}_{k\ge 1}$,
and denote the limit by $m$ and  $(X(0), Y(0),U,M)$, respectively.
By Skorokhod's theorem, on a probability space, 
there exist $\mathbb{R}^d\times\mathbb{R}^d\times C([0,1];\mathbb{R}^d)\times C([0,1];\mathbb{R}^d)$--valued random variables $(\tilde X_{k}(0), \tilde Y_{k}(0), \tilde U_{k},\tilde M_{k}), k\ge 1$ and
$(\tilde X(0), \tilde Y(0),\tilde U,\tilde M)$ such that
\begin{eqnarray*}
P^{(\tilde X_{k}(0), \tilde Y_{k}(0), \tilde U_{k},\tilde M_{k})}&=&
P^{(X_{n_k}(0), Y_{n_k}(0), U_{X_{n_k}},M_{X_{n_k}})},\quad k\ge 1\\
P^{(\tilde X(0), \tilde Y(0),\tilde U,\tilde M)}&=&P^{(X(0), Y(0),U,M)},\\
(\tilde X_{k}(0), \tilde Y_{k}(0), \tilde U_{k},\tilde M_{k}) &\to& (\tilde X(0), \tilde Y(0),\tilde U,\tilde M),\quad k\to\infty, \quad {\rm a.s..}
\end{eqnarray*}

Define $\{(\tilde X_{k},\tilde Y_{k})\}_{k\ge 1}$ by (\ref{3.8}) and (\ref{3.10})
with $(m,X(0), Y(0), U_{X},M_{X})$ replaced by $(m_k,\tilde X_{k}(0), \tilde Y_{k}(0), \tilde U_{k},\tilde M_{k})$.

If  $m_n\ge C, n\ge 1$, then 
$m\ne 0$ and $\{(\tilde X_{k},\tilde Y_{k})\}_{k\ge 1}$ is also convergent a.s. from Lemmas \ref{lm3.1}  and \ref{lm3.2}.
Indeed, 
\begin{eqnarray*}
&&\|\Psi^{m_{n_k}}(\tilde U_{k}+\tilde M_{k})-\Psi^{m}(\tilde U+\tilde M)\|_\infty\\
& \le& \|\Psi^{m_{n_k}}(\tilde U_{k}+\tilde M_{k}-(\tilde U+\tilde M))\|_\infty+
\|\Psi^{m_{n_k}}(\tilde U+\tilde M)-\Psi^{m}(\tilde U+\tilde M)\|_\infty\to 0,
\end{eqnarray*}
as $k\to\infty$,
from (\ref{3.2})--(\ref{3.3.0}).

Since $\Psi^{m_{n_k}}$ is Lipschitz continuous on $C([0,1];\mathbb{R}^d)$ from Lemma \ref{lm3.1},
the law of $(\tilde X_{k},\tilde Y_{k})$ is the same as $(X_{n_k}, Y_{n_k})$. 
Since  the space $C([0,1];\mathbb{R}^d)$ with the supnorm and $\mathbb{R}^d$ are Polish,
$\{Z_n\}_{n\ge 1}$ is tight. 

If $m_n\to 0$ as $n\to\infty$, then $m=0$ and
\begin{eqnarray}\label{3.18.1}
&&\|\Psi^{m_{n_k}}(\tilde U_{k}+\tilde M_{k})-(\tilde U+\tilde M)\|_\infty\\
& \le& \|\Psi^{m_{n_k}}(\tilde U_{k}+\tilde M_{k}-(\tilde U+\tilde M))\|_\infty
+
\|\Psi^{m_{n_k}}(\tilde U+\tilde M)-(\tilde U+\tilde M)\|_\infty\to 0,\nonumber
\end{eqnarray}
as $k\to\infty$, from (\ref{3.2}) and (\ref{3.4}).
In the same way as the case where $m\ne 0$,
(\ref{3.18.1}) together with (\ref{1.14}), (\ref{3.8}), and (\ref{3.10}) completes the proof.
\end{proof}

For $Z=(X,Y)\in \mathcal{A}^{m}$, let
\begin{equation}\label{3.19.0}
\overline X(t):=X(0)+\frac{1}{\gamma}(U_X(t)+M_X(t)),\quad 0\le t\le 1.
\end{equation}
The following lemma plays a crucial role  in the proofs of Theorem \ref{thm2.1} and Corollary \ref{co2.3}.

\begin{lemma}\label{lm3.4}
Suppose that (A1, i, ii) holds and that $\{m_n\}_{n\ge 1}$ is a sequence of positive real numbers 
that converges to $0$ as $n\to\infty$.
Then, for tight families $\{P_{0,n}\}_{n\ge 1}, \{P_{1,n}\}_{n\ge 1}\subset\mathcal{P}(\mathbb{R}^d)$,
 and $\{Z_n=(X_n, Y_n)\in\mathcal{A}^{m_n}(P_{0,n},P_{1,n})\}_{n\ge 1}$ such that (\ref{2.4}) holds
 and that $Y_n(0)\to 0$ as $n\to\infty$ weakly,
\begin{eqnarray}
\lim_{n\to\infty}\|Y_n\|_\infty&=&0,\quad {\rm in\hbox{ }Prob.},\label{3.19}\\
\lim_{n\to\infty}\|X_{n}-\overline X_{n}\|_\infty&=&0, \quad {\rm in\hbox{ }Prob.}.\label{3.20}
\end{eqnarray}
\end{lemma}

\begin{proof}
From Lemmas \ref{lm3.3}, $\{(X_{n},Y_{n}, U_{X_{n}},M_{X_{n}})\}_{n\ge 1}$ is tight.
Take a weakly convergent subsequence
$\{(X_{n_k},Y_{n_k}, U_{X_{n_k}},M_{X_{n_k}})\}_{k\ge 1}$.
In the same way as the proof of Lemmas \ref{lm3.3},
by Skorokhod's theorem, taking a different probability space if necessary, we assume that the convergence is almost sure. 
%\begin{equation}\label{3.21}
%\lim_{k\to\infty}|Y_{n_k}(0)|=0, \quad {\rm a.s.}.
%\end{equation}
%Indeed, from (\ref{3.1}) and (\ref{3.8}), in the same way as (\ref{3.14}),
%for $t\in (0,1]$,
%\begin{eqnarray*}
%|Y_{n_k}(0)|
%\le \gamma
%\left(1-\exp\left(-\frac{\gamma t}{m_{n_k}}\right)\right)^{-1}|X_{n_k}(t)-X_{n_k}(0)|
%+\sup_{0\le s\le t}|U_{X_{n_k}}(s)+M_{X_{n_k}}(s)|.
%\end{eqnarray*}
%Letting $k\to\infty$ and then $t\to 0$, (\ref{3.21})  holds.

From (\ref{3.4}) and (\ref{3.10}), 
\begin{eqnarray}\label{3.21.0}
&&\|Y_{n_k}\|_\infty\\
&\le&  \|\Psi^{m_{n_k}}(U_{X_{n_k}}+M_{X_{n_k}})-(U_{X_{n_k}}+M_{X_{n_k}})\|_\infty+|Y_{n_k}(0)|\to 0,\quad k\to\infty, {\rm a.s.},\nonumber
\end{eqnarray}
which also implies that the following holds:
\begin{equation}\label{3.22.0}
\lim_{k\to\infty}\|X_{n_k}-\overline X_{n_k}\|_\infty=0, \quad {\rm a.s.}.
\end{equation}
Indeed,  from (\ref{3.8}) and (\ref{3.10}),
\begin{eqnarray}\label{3.25.0}
X_{n}(t)-\overline X_{n}(t)
&=&\frac{1}{\gamma}\left\{\Psi^{m_{n}}(U_{X_{n}}+M_{X_{n}}+Y_{n}(0))(t)-(U_{X_{n}}(t)+M_{X_{n}}(t))\right\}\nonumber\\
\\
&=&\frac{1}{\gamma} (-Y_{n}(t)+Y_{n}(0)).\nonumber
\end{eqnarray}
Since (\ref{3.21.0})--(\ref{3.22.0}) hold for any weakly convergent subsequence and
since the convergence to $0$ a.s. implies that in probability, (\ref{3.19})--(\ref{3.20}) hold.
\end{proof}


The following lemma will be used in the proofs of Theorem \ref{thm2.3}
 and Corollary \ref{thm2.2}, and is given the proof for completeness
 (see (\ref{2.12})--(\ref{3.27}) for notation).

\begin{lemma}\label{lm3.5}
For $m>0$,
\begin{equation}
0\le \varphi^m(t)-t\le \frac{2m}{\gamma},\quad 0\le t\le 1.
\end{equation}
Equivalently,
\begin{equation}
0\le t-(\varphi^m)^{-1}(t)\le \frac{2m}{\gamma},\quad  \varphi^m(0)\le t\le 1.
\end{equation}
\end{lemma}

\begin{proof}
For $t\in [0,1]$,
\begin{eqnarray}\label{3.29}
\varphi^m(t)-t&=&\int_t^1 (1-f^m(s)^2)ds\ge 0.
\end{eqnarray}
The following together with (\ref{3.29}) completes the proof: for $s\in [0,1]$,
\begin{equation}
1-f^m(s)^2=2\exp\left(-\frac{\gamma (1-s)}{m}\right)-\exp\left(-\frac{2\gamma (1-s)}{m}\right)\le 2\exp\left(-\frac{\gamma (1-s)}{m}\right).
\end{equation}
\end{proof}

The following lemma plays a crucial role in the proof of Theorem \ref{thm2.3}.
\begin{lemma}\label{lm3.6}
Suppose that (A2, i) holds and that $m>0$.
Then  for any $P_0,P_1\in\mathcal{P}(\mathbb{R}^d)$, 
any $X\in\mathcal{A}(P_0,P_1)$, and for $Z^m=(X^m, Y^m)\in \mathcal{A}^{m}$ defined by (\ref{2.21})--(\ref{2.24}), the following holds:
for $t\in [0,1)$,
\begin{eqnarray}
X^m(t)&=&\left(1-\frac{K^m(t)}{K^m(1)}\right)X(0)
+f^m(t)\Psi^m\left(\frac{X(\varphi^m(\cdot))}{f^m(\cdot)^2}\right)(t),\label{3.26}\\
Y^m(t)&=&-\frac{\exp(-\gamma t/m)}{K^m(1)}X(0)
+\frac{X(\varphi^m(t))}{K^m(1-t)}
-\gamma\Psi^m\left(\frac{X(\varphi^m(\cdot))}{f^m(\cdot)^2}\right)(t).\qquad\label{2.27}
\end{eqnarray}
\end{lemma}

\begin{proof}

From (\ref{2.21})--(\ref{2.24}), and (\ref{3.9}), 
\begin{eqnarray}
X^m(t)&=&X(0)+K^m(t)\frac{X(\varphi^m(0))-X(0)}{K^m(1)}\label{4.13.0}\\
&&\qquad +\int_{0}^{t} K^m(t-s)(u_X^m(s)f^m(s)ds +d W^m(s)).\nonumber
\end{eqnarray}

From (\ref{4.13.0}),
%--(\ref{2.24}) and (\ref{3.9}), 
 the following implies (\ref{3.26}) (see (\ref{3.7})--(\ref{2.12})
for notation): for $t\in [0,1)$,
\begin{eqnarray}\label{3.40.0}
&&\int_{0}^{t} K^m(t-s)(u_X^m(s)f^m(s)ds +d W^m(s))\\
&=&-\frac{K^m(t)}{K^m(1)}X(\varphi^m(0))
+f^m(t)\int_{0}^{t}
\frac{\gamma}{m}\exp\left(-\frac{\gamma (t-s)}{m}\right)\frac{X(\varphi^m(s))}{f^m(s)^2}ds.\nonumber
 \end{eqnarray}
We prove (\ref{3.40.0}).  
 From (\ref{3.23}), by the integration by parts,
\begin{eqnarray}
&&\int_{0}^{t} K^m(t-s)(u_X^m(s)f^m(s)ds +d W^m(s))\label{4.15.0}\\
&=&\int_{0}^{t} \frac{K^m(t-s)}{K^m(1-s)}dX(\varphi^m(s))\nonumber\\
 &=&\left[\frac{K^m(t-s)}{K^m(1-s)}X(\varphi^m(s))\right]_{s=0}^t
 -\int_{0}^{t} \frac{d}{ds}\left\{\frac{K^m(t-s)}{K^m(1-s)}\right\}X(\varphi^m(s))ds\nonumber\\
 &=&-\frac{K^m(t)}{K^m(1)}X(\varphi^m(0))
 -\int_{0}^{t} \frac{d}{ds}\left\{\frac{K^m(t-s)}{K^m(1-s)}\right\}X(\varphi^m(s))ds,\nonumber
 \end{eqnarray}
 \begin{eqnarray}\label{3.39}
&& -\frac{d}{ds}\left\{\frac{K^m(t-s)}{K^m(1-s)}\right\}\\
&=&\frac{1}{m}\left\{\frac{\exp(-\gamma (t-s)/m)K^m(1-s)-K^m(t-s)\exp(-\gamma (1-s)/m)}{K^m(1-s)^2}\right\},\nonumber
\end{eqnarray}
\begin{eqnarray}\label{3.40}
&&\exp\left(-\frac{\gamma (t-s)}{m}\right)K^m(1-s)-K^m(t-s)\exp\left(-\frac{\gamma (1-s)}{m}\right)\\
%&=&\exp\left(-\frac{\gamma (t-s)}{m}\right)\frac{1}{\gamma}\left(1-\exp\left(-\frac{\gamma (1-s)}{m}\right)\right)\nonumber\\
%&&\qquad -\frac{1}{\gamma}\left(1-\exp\left(-\frac{\gamma (t-s)}{m}\right)\right)\exp\left(-\frac{\gamma (1-s)}{m}\right)\nonumber\\
%&=&\frac{1}{\gamma}\left\{\exp\left(-\frac{\gamma (t-s)}{m}\right)-\exp\left(-\frac{\gamma (1-s)}{m}\right)\right\}\nonumber\\
&=&\exp\left(-\frac{\gamma (t-s)}{m}\right)K^m(1-t).\nonumber
\end{eqnarray}

From (\ref{2.21})--(\ref{2.24}) and (\ref{3.11}), the following implies (\ref{2.27}) (see (\ref{3.7})--(\ref{2.12})
for notation): for $t\in [0,1)$,
\begin{eqnarray}
&&\int_{0}^{t}\exp\left(-\frac{\gamma (t-s)}{m}\right)(u_X^m(s)f^m(s)ds +d W^m(s))\label{4.23}\\
 &=&\frac{X(\varphi^m(t))}{K^m(1-t)}-\frac{\exp(-\gamma t/m)}{K^m(1)}X(\varphi^m(0))
-\gamma\Psi^m\left(\frac{X(\varphi^m (\cdot))}{f^m(\cdot)^2}\right)(t).\nonumber
\end{eqnarray}
Indeed,
\begin{eqnarray}
Y^m(t)&=&\exp\left(-\frac{\gamma t}{m}\right)\frac{X(\varphi^m(0))-X(0)}{K^m(1)}\label{4.19}\\
&&\qquad +\int_{0}^{t} \exp\left(-\frac{\gamma (t-s)}{m}\right)(u_X^m(s)f^m(s)ds +d W^m(s)).\nonumber
\end{eqnarray}

We prove (\ref{4.23}).
From   (\ref{3.23}), in the same way as (\ref{4.15.0}), by the integration by parts,
\begin{eqnarray}
&&\int_{0}^{t}\exp\left(-\frac{\gamma (t-s)}{m}\right)(u_X^m(s)f^m(s)ds +d W^m(s))\\%\label{4.23}\\
%&=&\int_{0}^{t} \frac{\exp(-\gamma (t-s)/m)}{K^m(1-s)}dX(\varphi^m(s))\nonumber\\
&=&\left[\frac{\exp(-\gamma (t-s)/m)}{K^m(1-s)}X(\varphi^m(s))\right]_{s=0}^t
 -\int_{0}^{t} \frac{d}{ds}\left\{\frac{\exp(-\gamma (t-s)/m)}{K^m(1-s)}\right\}X(\varphi^m(s))ds \nonumber\\
 &=&\frac{X(\varphi^m(t))}{K^m(1-t)}-\frac{\exp(-\gamma t/m)}{K^m(1)}X(\varphi^m(0))
-\gamma\Psi^m\left(\frac{X(\varphi^m (\cdot))}{f^m(\cdot)^2}\right)(t)\nonumber
 %&=&\frac{1}{K^m(1-t)}X(\varphi^m(t))-\frac{\exp(-\gamma t/m)}{K^m(1)}X(\varph^m(0))
% -\int_{0}^{t} \frac{d}{ds}\left\{\frac{K^m(t-s)}{K^m(1-s)}\right\}X(\varphi^m(s))ds,\nonumber
\end{eqnarray}
from (\ref{3.39})--(\ref{3.40}).
Indeed, from (\ref{3.7}),
$$\exp\left(-\frac{\gamma (t-s)}{m}\right)=-\gamma K^m(t-s)+1,\quad 0\le s\le t\le 1,$$
$$\frac{d}{ds}\left\{\frac{1}{K^m(1-s)}\right\}=\frac{1}{mK^m(1-s)^2}
\exp\left(-\frac{\gamma (t-s)}{m}\right)\exp\left(-\frac{\gamma (1-t)}{m}\right).
$$

\end{proof}

























%\newpage
\section{Proofs of main results}\label{sec:4}

In this section, we prove our results.

For $m>0$, $t\in [0,1], z=(x,y)%, z_1=(x_1,y_1)
\in  \mathbb{R}^{d}\times \mathbb{R}^{d}, u\in \mathbb{R}^{d}$, and $f\in C^{1,2}([0,1]\times \mathbb{R}^{2d})$,
let 
\begin{eqnarray}
\mathcal{L}^m_{t,z,u}f(t,z)
&:=&\frac{\partial }{\partial t}f(t, z)+\langle D_x f(t, z),\frac{1}{m}y\rangle+\langle D_y f(t, z),u-\frac{\gamma}{m}y\rangle\\
&&\qquad +\frac{1}{2}Trace (a(t,z)D_y^2f(t, z))\nonumber
\end{eqnarray}
(see (\ref{2.56}) for notation), where $D_x:=(\partial/\partial x_i)_{i=1}^d$ and $D_y^2:=(\partial/\partial y_i\partial y_j)_{i,j=1}^d$.

Modifying the idea in \cite{M08,M21} (see also \cite{M2021} and the reference therein), 
we prove Theorem \ref{pp1.1} by Bogachev--R\"ockner--Shaposhnikov's superposition principle (see \cite{Bog21}).

\begin{proof}[Proof of Theorem \ref{pp1.1}]
From Lemma \ref{lm3.3}, $\{Q_n\}_{n\ge 1}$ and $\{Z_n\}_{n\ge 1}$ are tight.
Take weakly convergent subsequences $\{Q_{n_k}\}_{k\ge 1}$ and $\{Z_{n_k}\}_{k\ge 1}$ of $\{Q_n\}_{n\ge 1}$
 and $\{Z_n\}_{n\ge 1}$, respectively.
Let $Q_\infty$ and $Z_\infty$ denote the limits of $Q_{n_k} $ and $\{Z_{n_k}\}_{k\ge 1}$
as $k\to\infty$, respectively.
Then 
\begin{equation}\label{4.2}
Q_\infty(dt\hbox{ }dz\times \mathbb{R}^d)=dtP^{Z_\infty(t)}(dz),
\end{equation}
since 
$$\int_{[0,1]\times \mathbb{R}^{2d}}f(t,z)Q_{n_k}(dt\hbox{ }dz\times \mathbb{R}^d)
=\int_0^1 E[f(t,Z_{n_k}(t))]dt,\quad f\in C_b([0,1]\times \mathbb{R}^{2d}).
$$

We prove (\ref{2.5})--(\ref{2.6}).
For any $f\in C_0^{2}(\mathbb{R}^{2d})$ and $t\in [0,1]$, from (\ref{4.2}),
\begin{eqnarray}\label{4.3.0}
&&\int_{ \mathbb{R}^{2d}}f(z)P^{Z_\infty(t)}(dz)-
\int_{ \mathbb{R}^{2d}}f(z)P^{Z_\infty(0)}(dz)\\
&=&\int_{[0,t]\times \mathbb{R}^{2d}}
\left\{\mathcal{L}^m_{s,z, E^{Q_\infty}[u|s,z]}f(z)\right\}dsP^{Z_\infty(s)}(dz).\nonumber
\end{eqnarray}
Indeed, for any $\varphi\in C_0^{1,2}([0,1]\times \mathbb{R}^{2d})$, by the It\^o formula,  
\begin{eqnarray}\label{4.4}
0&=&\int_{[0,1]}dt
E\left[\mathcal{L}^m_{t,Z_{n_k}(t),u_{X_{n_k}}(t)}\varphi(t,Z_{n_k}(t))\right]\\
&=&\int_{[0,1]\times \mathbb{R}^{2d}\times \mathbb{R}^{d}}
\mathcal{L}^m_{t,z,u}\varphi(t,z)Q_{n_k} (dt\hbox{ }dz\hbox{ }du)\nonumber\\
&\to&\int_{[0,1]\times \mathbb{R}^{2d}\times \mathbb{R}^{d}}
\mathcal{L}^m_{t,z,u}\varphi(t,z)Q_\infty(dt\hbox{ }dz\hbox{ }du),
\quad k\to\infty\nonumber\\
&=&\int_{[0,1]\times \mathbb{R}^{2d}}
\mathcal{L}^m_{t,z, E^{Q_\infty}[u|t,z]}\varphi(t,z)Q_\infty(dt\hbox{ }dz\times \mathbb{R}^d),\nonumber
\end{eqnarray}
from (A1, i, ii).
Here we used the following: from (\ref{2.4}), for $R>0$,
\begin{eqnarray}\label{4.5.0}
&&\int_{[0,1]\times \mathbb{R}^{2d}\times B_{d,R}^c} |u|Q_{n_k} (dt\hbox{ }dz\hbox{ }du)\\
& \le& \frac{1}{C_{1,R}
}\int_{[0,1]\times \mathbb{R}^{2d}\times B_{d,R}^c} L(t,z;u)Q_{n_k} (dt\hbox{ }dz\hbox{ }du)\nonumber\\
&\le&  \frac{1}{C_{1,R}
}
E\left[\int_0^1 L(t,Z_{n_k}(t);u_{X_{n_k}}(t))dt\right] 
\to 0,\nonumber
\end{eqnarray}
as $R\to\infty$, uniformly in $k$ (see the proof of Lemma \ref{lm3.3} for notation).


%There exists $\{q(t,dz)\}_{0\le t\le 1}\subset \mathcal{P}(\mathbb{R}^{2d})$
%such that
%\begin{equation}\label{4.5}
%Q_\infty(dt\hbox{ }dz\times \mathbb{R}^d)=dt q(t,dz)
%\end{equation}
%since 
%\begin{equation*}
%Q_{\infty} (dt\times \mathbb{R}^{2d}\times \mathbb{R}^{d})=
%\lim_{k\to\infty}
%Q_{n_k} (dt\times \mathbb{R}^{2d}\times \mathbb{R}^{d})
%=dt.
%\end{equation*}
%Notice that we can take a weakly continuous version of $t\mapsto q(t,dz)$
%(see e.g., Remark 2.3 in \cite{1-Tre}).
%From (\ref{4.4}) and (\ref{4.5}), 

From (\ref{4.3.0}), Bogachev--R\"ockner--Shaposhnikov's superposition principle implies that 
there exists $Z=(X,Y)\in \mathcal{A}^m$ such that 
\begin{equation}\label{4.6.0}
P^{Z(t)}(dz)=P^{Z_\infty(t)}(dz),\quad 0\le t\le 1,
\end{equation}
and that (\ref{2.5})--(\ref{2.6}) hold (see \cite{Bog21}, Theorem 1.1).
Here notice that from (\ref{4.5.0}), for $R>0$, 
\begin{eqnarray*}
&&\int_{[0,1]\times \mathbb{R}^{2d}}
|E^{Q_\infty}[u|t,z]|Q_\infty(dt\hbox{ }dz\times \mathbb{R}^d)\nonumber\\
&\le& \int_{[0,1]\times \mathbb{R}^{2d}\times \mathbb{R}^{d}} |u|Q_{\infty} (dt\hbox{ }dz\hbox{ }du)\\
& \le& R+\liminf_{k\to\infty}\frac{1}{C_{1,R}
}\int_{[0,1]\times \mathbb{R}^{2d}\times B_{d,R}^c} L(t,z;u)Q_{n_k} (dt\hbox{ }dz\hbox{ }du)<\infty.
\end{eqnarray*}
$P^{X(t)}=P_t, t=0,1$ and $P^{Y(0)}\in B $ from (\ref{4.6.0}), since $Z_{n_k}\in\mathcal{A}^m (B ,P_0;P_1)$
and since $B$ is a closed set.

We prove (\ref{1.9.1}).
Taking a subsequence if necessary, we assume that the following  is convergent:
$$E\biggl[\int_0^1 L(t,Z_n(t);u_{X_n} (t))dt \biggr].$$
Take a weakly convergent subsequence $\{Q_{n_k}\}_{k\ge 1}$ of $\{Q_{n}\}_{n\ge 1}$,
its weak limit $Q_\infty$, and $Z=(X,Y)\in \mathcal{A}^m(B,P_0;P_1)$ that satisfies (\ref{2.5})--(\ref{2.6}).
By Skorokhod's theorem, on a probability space, there exist $[0,1]\times \mathbb{R}^{2d}\times \mathbb{R}^{d}$--valued random variables $(T_k, Z_k, U_k), k\ge 1$ and
$(T, Z, U)$ such that 
$$P^{(T_k, Z_k, U_k)}=Q_{n_k},\quad k\ge 1,\quad P^{(T, Z, U)}=Q_\infty,$$
$$\lim_{k\to\infty} (T_k, Z_k, U_k)=(T, Z, U),\quad {\rm a.s.}.$$
By Fatou's lemma and Jensen's inequality, from (A1, iii, iv), the following holds:
\begin{eqnarray}\label{4.7.0}
&&\lim_{k\to\infty}E\biggl[\int_0^1 L(t,Z_{n_k}(t);u_{X_{n_k}} (t))dt \biggr]\\
&=&\lim_{k\to\infty}
\int_{[0,1]\times \mathbb{R}^{2d}\times \mathbb{R}^{d}}
L(t,z;u)Q_{n_k}(dt\hbox{ }dz\hbox{ }du)
=\lim_{k\to\infty}E[L(T_k, Z_k;U_k)]\nonumber\\
& \ge &E[L(T, Z; U)]=\int_{[0,1]\times \mathbb{R}^{2d}\times \mathbb{R}^{d}}
L(t,z;u)Q_\infty (dt\hbox{ }dz\hbox{ }du)\nonumber\\
& \ge&\int_{[0,1]\times \mathbb{R}^{2d}}
L(t,z;E^{Q_\infty}[u|t,z])Q_\infty (dt\hbox{ }dz\times \mathbb{R}^d)
=E\biggl[\int_0^1 L(t,Z(t);u_{X} (t))dt \biggr].\nonumber
\end{eqnarray}

\end{proof}


Theorem \ref{thm2.1} can be proven in the same way as Theorem \ref{pp1.1}.
The key is Lemma \ref{lm3.4} which implies (\ref{4.7}) given below.


\begin{proof}[Proof of Theorem \ref{thm2.1}]
$\{Q_n\}_{n\ge 1}$ and $\{Z_n=(X_n, Y_n)\}_{n\ge 1}$ are tight,
from Lemma \ref{lm3.3} and  from assumption.
%$$P^{Y_\infty(0)}\in B_n, \quad n\ge 1$$
%since $B_n$ is closed and 
%$$P^{Y_{n_k}(0)}\in B_{n_k}\subset B_n, \quad n_k\ge n.$$
%In particular, 


Take weakly convergent subsequences 
$\{Q_{n_k}\}_{k\ge 1}$ and 
$\{Z_{n_k}\}_{k\ge 1}$
of $\{Q_n \}_{n\ge 1}$ and $\{Z_n\}_{n\ge 1}$, respectively, and let $\overline Q_\infty$ denote the limit of $Q_{n_k} $ as $k\to\infty$.
Then from Lemma \ref{lm3.4},
\begin{equation}
\overline Q_\infty(dt\hbox{ }dz\hbox{ }du)=dt\delta_0(dy)\overline Q_\infty(dt\hbox{ }dx\times \mathbb{R}^d \times du).
\end{equation}
We can prove (\ref{2.13})--(\ref{2.14}) in the same way as Theorem \ref{pp1.1},
from Lemma \ref{lm3.4} and from the following:
for any $f\in C_0^{1,2}([0,1]\times \mathbb{R}^{d})$, 
\begin{eqnarray}\label{4.7}
0&=&\int_{[0,1]\times \mathbb{R}^{2d}} \left\{
\frac{\partial }{\partial t}f(t, x)+
\frac{1}{\gamma}\langle D_x f(t, x),u\rangle\right.\\
&&\qquad \left.
+\frac{1}{2\gamma^2}Trace (a(t, z)D_x^2f(t, x))\right\}
\overline Q_\infty(dt\hbox{ }dz\hbox{ }du)\nonumber
\end{eqnarray}
(see (\ref{4.5.0})).
Indeed,  by the It\^o formula, 
\begin{eqnarray}%\label{4.4}
0&=&E\left[\int_0^1 \left\{
\frac{\partial }{\partial t}f(t, \overline X_{n_k}(t))+
\frac{1}{\gamma}\langle D_x f(t, \overline X_{n_k}(t)),u_{X_{n_k}(t)}\rangle\right.\right.\nonumber\\
&&\qquad\left.\left.+\frac{1}{2\gamma^2}Trace (a(t, Z_{n_k}(t))D_x^2f(t, \overline X_{n_k}(t))\right\}dt\right]\nonumber
\end{eqnarray}
(see (\ref{3.19.0})).
From Lemma \ref{lm3.4}, %by the bounded convergence theorem
\begin{eqnarray}\label{4.10}
&&E\left[\int_0^1 \left\{
\left|\frac{\partial }{\partial t}f(t, \overline X_{n_k}(t)) -\frac{\partial }{\partial t}f(t,X_{n_k}(t))\right|
\right.\right.\\
&&\qquad \left.\left.
+\left|D_x^2f(t, \overline X_{n_k}(t))-D_x^2f(t, X_{n_k}(t))\right|\right\}
dt\right]\to 0,\quad k\to\infty.\nonumber
\end{eqnarray}
From (A1, ii), for $R>0$,
\begin{eqnarray}\label{4.11}
&&E\left[\int_0^1 
|D_x f(t, \overline X_{n_k}(t))-D_x f(t, X_{n_k}(t))||u_{X_{n_k}(t)}|dt\right]\\
& \le& R \cdot E\left[\int_0^1 |D_x f(t, \overline X_{n_k}(t))-D_x f(t, X_{n_k}(t))|dt\right]\nonumber\\
&&\qquad +2\sup_{(s,x)\in [0,1]\times \mathbb{R}^d}|D_x f(s,x)|\frac{1}{C_{1,R}
}
E\left[\int_0^1 L(t,Z_{n_k}(t);u_{X_{n_k}(t)})dt\right]
\to 0,\nonumber
\end{eqnarray}
as $k\to\infty$ and then $R\to\infty$.

In the same way as (\ref{4.7.0}), we can prove (\ref{2.15}).
\end{proof}

Lemma \ref{lm3.2} plays a crucial role in the proof of Theorem \ref{thm2.3}.
The key idea that leads to $(X^m, Y^m)\in\mathcal{A}^m(P_0,P_1)$ to approximate 
$X\in\mathcal{A}(P_0,P_1)$ is the following.
For $(\hat X, \hat Y)\in\mathcal{A}^m(P_0,P_1)$,
substitute $t=1$ in (\ref{3.9}). Then 
$$\hat X(1)=\hat X(0)+K^m(1) \hat Y(0)+\int_{0}^{1} K^m(1-s)(u_{\hat X}(s)ds +d W(s)),
$$
which implies that $\hat X(t)$ is close to the following semimartingale:
$$\xi (t)=\hat X(0)+K^m(1) Y(0)+\int_{0}^{t} K^m(1-s)(u_{\hat X}(s)ds +d W(s)).$$
(\ref{2.24.00}) implies the we should consider the time change of $X$ under (A2, i).

\begin{proof}[Proof of Theorem \ref{thm2.3}]

From (\ref{3.23}) and (\ref{4.13.0}), 
\begin{equation}
\quad X^m(1)%&=&X(0)+K^m(1) Y^m(0)+\int_{0}^{1} K^m(1-s)(u_X^m(s)f^m(s)ds +d W^m(s))\\
=X(\varphi^m(0))+\int_{0}^{1} K^m(1-s)(u_X^m(s)f^m(s)ds +d W^m(s))=X(1)
\end{equation}
since $\varphi^m(1)=1$ (see  (\ref{3.7}) and (\ref{3.27}) for notation).

We prove that $X^m$ converges  to $X$, as $m\to 0$, locally uniformly on $[0,1)$, a.s..
From (\ref{3.26}), for $t\in [0,1)$,
\begin{eqnarray}%\label{4.13.0}
&&X^m(t)-X(t)\\
&=&\left(1-\frac{K^m(t)}{K^m(1)}\right)(X(0)-X(t))
-\left(\frac{1}{K^m(1)}-\frac{1}{K^m(1-t)}\right)K^m(t)X(t)\nonumber\\
&&\qquad -\frac{K^m(t)}{K^m(1-t)}(X(t)-X(\varphi^m(t)))\nonumber\\
&&\qquad 
+f^m(t)\left\{\Psi^m\left(\frac{X(\varphi^m(\cdot))}{f^m(\cdot)^2}\right)(t)
-\gamma K^m(t)\frac{X(\varphi^m(t))}{f^m(t)^2}\right\}\nonumber
\end{eqnarray}
since $\gamma K^m (1-t)=f^m(t)$ (see (\ref{2.12}) for notation).

For $\delta\in (0,1)$ and $t\in [0,1]$, considering the cases when $t\in [0,\delta]$ and when  $t\in [\delta,1]$,
\begin{eqnarray}
&&\left|\left(1-\frac{K^m(t)}{K^m(1)}\right)(X(0)-X(t))\right|\\
& \le& \sup_{0\le s\le \delta}|X(0)-X(s)|
+\frac{\exp(-\gamma\delta/m)}{1-\exp(-\gamma/m)}\times 2\|X\|_{\infty}\to 0,\nonumber
\end{eqnarray}
as $m\to 0$ and then $\delta\to 0$, since 
\begin{equation}\label{4.21}
0\le K^m(t)=\frac{1-\exp(-\gamma t/m)}{\gamma}\le K^m(1), \quad 0\le t\le 1.
\end{equation}

Take $t_0\in (0,1)$. For $t\in [0, t_0]$,
\begin{equation}%\label{3.26}
\left|\left(\frac{1}{K^m(1)}-\frac{1}{K^m(1-t)}\right)K^m(t)X(t)\right|
\le \frac{\exp(-\gamma(1-t_0)/m)}{1-\exp(-\gamma(1-t_0)/m)}\|X\|_\infty\to 0,
\end{equation}
as $m\to 0$,  from (\ref{4.21}), since $t\mapsto K^m(1-t)$ is decreasing.

From Lemma \ref{lm3.5}, 
%for sufficiently small $m>0$, the following holds: 
%\begin{eqnarray}\label{3.46}
%0\le t& \le \varphi^m(t)\le t+\frac{2m}{\gamma}\le 1,\quad t\in[0,t_0].
%\end{eqnarray}
 for $t\in [0, t_0]$,
\begin{eqnarray}%\label{3.26}
&&\left|\frac{K^m(t)}{K^m(1-t)}(X(t)-X(\varphi^m(t)))\right|\\
& \le& \frac{1}{1-\exp(-\gamma(1-t_0)/m)}\sup_{0\le s_1\le s_2\le 1, |s_1-s_2|\le 2m/\gamma}|X(s_1)-X(s_2)|\to 0, \nonumber
\end{eqnarray}
as $m\to 0$, from (\ref{4.21}), since $t\mapsto K^m(1-t)$ is decreasing.

For $t\in [0, t_0]$ and $\delta\in (0,1)$, from (\ref{3.3}) and Lemma \ref{lm3.5},
\begin{eqnarray}
&&\left|\Psi^m\left(\frac{X(\varphi^m(\cdot))}{f^m(\cdot)^2}\right)(t)
-\gamma K^m(t)\frac{X(\varphi^m(t))}{f^m(t)^2}\right|\label{3.34}\\
& \le& 2\left\|\frac{X(\varphi^m(\cdot))}{f^m(\cdot)^2}\right\|_{\infty,t_0}\exp\left (-\frac{\gamma\delta}{m} \right)
+\sup_{t-t\wedge\delta\le s\le t}\left|\frac{1}{f^m(s)^2}-\frac{1}{f^m(t)^2}\right|
| X(\varphi^m(s))|\nonumber\\
&&\qquad +\sup_{t-t\wedge\delta\le s\le t}\left|\frac{X(\varphi^m(s))-X(\varphi^m(t))}{f^m(t)^2}\right|\nonumber\\
& \le& 2\|X\|_\infty\left(\frac{1}{f^m(t_0)^2}\exp\left (-\frac{\gamma\delta}{m} \right)
+\frac{1}{f^m(t_0)^4}\exp\left(-\frac{\gamma(1-t_0)}{m}\right)\right)\nonumber\\
&&\qquad +\frac{1}{f^m(t_0)^2}\sup_{0\le s_1\le s_2\le 1, |s_1-s_2|\le \delta+2m/\gamma}
|X(s_1)-X(s_2)|\to 0,\nonumber
\end{eqnarray}
as $m\to 0$ and then $\delta\to 0$.
Indeed,  $t\mapsto f^m (t)=1-\exp (-\gamma (1-t)/m)$ is decreasing and 
$$\left|\frac{1}{f^m(s)^2}-\frac{1}{f^m(t)^2}\right|\le \frac{1}{f^m(t)^4}\times 2\exp\left(-\frac{\gamma (1-t)}{m}\right),\quad 0\le s\le t\le 1.$$



We   prove that
$Y^m$ converges  to $0$, as $m\to 0$, locally uniformly on $[0,1)$, a.s..
Take $t_0\in (0,1)$. %  and sufficiently small $m>0$ such that (\ref{3.46}) holds.
Then, from (\ref{2.27}) and (\ref{3.34}), we only have to prove, to complete the proof,  that the following converges to $0$ uniformly in $t\in [0, t_0]$, as $m\to 0$:
\begin{eqnarray}\label{4.24}
&&-\frac{\exp(-\gamma t/m)}{K^m(1)}X(0)+
\frac{X(\varphi^m(t))}{K^m(1-t)}-\gamma^2K^m(t)\frac{X(\varphi^m (t))}{f^m(t)^2}\\
%& =\left|\frac{K^m(1-t)-K^m(t)}{K^m(1-t)^2}X(\varphi^m(t))-\frac{\exp(-\gamma t/m)}{K^m(1)}X(\varphi^m(0))\right|\nonumber\\
&=  &\frac{\exp(-\gamma t/m)}{K^m(1)}(X(\varphi^m(t))-X(0))\nonumber\\
&&\qquad +
\frac{(K^m(1-t)-K^m(t))K^m(1)-K^m(1-t)^2\exp(-\gamma t/m)}{K^m(1-t)^2K^m(1)}X(\varphi^m(t)),\nonumber
\end{eqnarray}
from (\ref{2.12}). 
We prove that (\ref{4.24}) converges to $0$ uniformly in $t\in [0, t_0]$, as $m\to 0$, a.s..
For $t\in [0, t_0]$ and $\delta\in (0,t_0)$, from Lemma \ref{lm3.5},
considering the cases when $t\in [0,\delta]$ and when  $t\in [\delta,t_0]$,
\begin{eqnarray}
&&\frac{\exp(-\gamma t/m)}{K^m(1)}|X(\varphi^m(t))-X(0)|\\
&\le&  \frac{1}{K^m(1)}\sup_{0\le s\le\min(\delta+2m/\gamma,1)} |X(s)-X(0)|
+\frac{\exp(-\gamma \delta/m)}{K^m(1)}\times 2\|X\|_\infty
\to 0,\nonumber
\end{eqnarray}
as $m\to 0$ and then $\delta\to 0$.
\begin{eqnarray}
&&\left|\frac{(K^m(1-t)-K^m(t))K^m(1)-K^m(1-t)^2\exp(-\gamma t/m)}{K^m(1-t)^2K^m(1)}X(\varphi^m(t))\right|\qquad\quad\\
&\le& 
\frac{\exp(-\gamma (1-t_0)/m)}{\gamma^2K^m(1-t_0)^3}\|X\|_\infty \to 0,\nonumber
\end{eqnarray}
as $m\to 0$,  since
\begin{eqnarray*}
&&\gamma^2|(K^m(1-t)-K^m(t))K^m(1)-K^m(1-t)^2\exp(-\gamma t/m)|\\
&=&\exp\left(-\frac{\gamma (1-t)}{m}\right)+\exp\left(-\frac{\gamma (1+t)}{m}\right)
-2\exp\left(-\frac{\gamma }{m}\right)\le \exp\left(-\frac{\gamma (1-t)}{m}\right),
\end{eqnarray*}
and since $t\mapsto K^m(1-t)$ is decreasing. 
Here, notice that $x\mapsto \exp (-x)$ is convex.

We prove (\ref{2.29}).
From (A2, iii),
\begin{eqnarray}
&&E\left[\int_0^1L(t, Z^m(t);u_X^m(t)f^m(t))dt\right]\label{4.28}\\
& \le& 
E\left[\int_0^1R_1(t, Z^m(t); u_X(\varphi^m(t)))f^m(t)^2dt\right]
+E\left[\int_0^1L(t, Z^m(t); 0)dt\right]\nonumber\\
&=&
E\left[\int_{\varphi^m(0)}^1L((\varphi^m)^{-1}(s), Z^m((\varphi^m)^{-1}(s));u_X(s))ds\right]\nonumber\\
&&\qquad +E\left[\int_0^1(-f^m(t)^2+1)L(t, Z^m(t); 0)dt\right].\nonumber
\end{eqnarray}

From Lemma \ref{lm3.5}, for $s\in [\varphi^m(0),1]$,
\begin{eqnarray}\label{4.13}
&&L((\varphi^m)^{-1}(s), Z^m((\varphi^m)^{-1}(s));u_X(s))
\\
& \le& \Delta L\left(\frac{2m}{\gamma},\infty\right)
(1+L(s, X(s), 0; u_X(s)))+L(s, X(s), 0;u_X(s)).\nonumber
\end{eqnarray}
From (A2, ii, iv), $L(t, Z^m(t); 0)$ is bounded (see Remark \ref{rk2.1}, (iii)).
(\ref{2.29}) holds from (A2, iv).

We prove (\ref{2.30}).
From Lemmas \ref{lm3.5} and what we have proven above, 
$$((\varphi^m)^{-1}(s), X^m((\varphi^m)^{-1}(s)), Y^m((\varphi^m)^{-1}(s)))\to (s, X(s), 0), \quad m\to 0,$$
locally uniformly in $s\in (0,1)$, a.s..
$\varphi^m(0)\to 0$, as $m\to 0$.
$L(t, Z^m(t); 0)$ is bounded, and 
\begin{eqnarray*}
0&\le& -f^m(t)^2+1\le 1,  \quad 0\le t<1,\\
1-f^m(t)^2&\to& 0, \quad m\to 0, \quad 0\le t<1.
\end{eqnarray*}
(A2, ii) and (\ref{4.13}) complete the proof by Lebesgue's dominated convergence theorem.

\end{proof}




We prove Corollary \ref{thm2.2}
by showing that $P^{Y^m(0)}\in \mathcal{P}_{1,C(m,P_0,P_1)}(\mathbb{R}^d)$ for $X\in\mathcal{A}(P_0;P_1)$ for which $E\left[\int_0^1L_0(t,X(t);u_X(t))dt\right]$ is sufficiently small.
 
\begin{proof}[Proof of Corollary \ref{thm2.2}]
Take $X\in\mathcal{A}(P_0;P_1)$ such that 
$$E\left[\int_0^1L_0(t,X(t);u_X(t))dt\right]\le V^0(P_0,P_1)+1.$$
Then,  
\begin{equation}
E[|X(\varphi^m (0))-X(0)|]\le  K^m(1)C(m,P_0,P_1), \quad m>0,
\end{equation}
and $Z^m=(X^m, Y^m)\in \mathcal{A}^m(P_0, P_1)$ defined by  (\ref{2.21})--(\ref{2.24})
belongs to $\mathcal{A}^m(\mathcal{P}_{1,C(m,P_0,P_1)}(\mathbb{R}^d),P_0;P_1)$ for $m>0$.
Indeed, for $R>0$,
\begin{eqnarray*}
&&E[|X(\varphi^m (0))-X(0)|]\\
&\le& \frac{1}{\gamma}\left\{\int_0^{\varphi^m (0)}E[|u_X(t)|]dt+E[|W(\varphi^m (0))|]\right\}\\
&\le& \frac{1}{\gamma}\left\{\varphi^m (0) R+\frac{1}{C_{1,R}}E
\left[\int_0^1L_0(t,X(t);u_X(t))dt\right]+E[|W(\varphi^m (0))|]\right\}.
\end{eqnarray*}

Since $\mathcal{P}_{1,C(m,P_0,P_1)}(\mathbb{R}^d)\subset D_m$, for $m\in (0,
\varepsilon_0\gamma/2]$, from (\ref{2.29}), 
\begin{eqnarray}
V^m(D_m,P_0;P_1)&\le &V^m(\mathcal{P}_{1,C(m,P_0,P_1)}(\mathbb{R}^d),P_0;P_1)\\
&\le&
\sup_{\overline m\in (0, \varepsilon_0\gamma/2]}E\left[\int_0^1L(t, Z^{\overline m}(t);u_X^{\overline m}(t)f^{\overline m}(t))dt\right]<\infty.\nonumber
\end{eqnarray}
(\ref{2.30}) completes the proof.
\end{proof}

In the proof of Theorem\ref{thm2.6},
we show that  $(\mathcal{X}^m,\mathcal{Y}^m)$ defined by (\ref{2.42})--(\ref{2.43})
is close to $(X^m, Y^m )$ defined by (\ref{2.21})--(\ref{2.24}) as $m\to 0$ and make use of 
Theorem \ref{thm2.3}.

\begin{proof}[Proof of Theorem\ref{thm2.6}]
From  (\ref{2.42})--(\ref{2.43}), (\ref{3.9}), (\ref{3.11}), (\ref{4.13.0}), and (\ref{4.19}), 
\begin{eqnarray}
\mathcal{X}^m(t)
 &=&X(0)+K^m(t)\mathcal{Y}^m(0)\label{4.35}\\
 &&\qquad +\int_0^tK^m(t-s)
\left(\left(u_{X}^m(s)
+\beta^m(0)\right)f^m(s)ds+dW^m(s)\right)\nonumber\\
&=&X^m(t)-K^m(t)\frac{X(\varphi^m(0))-X(0)-K^m(1)\mathcal{Y}^m(0)}{K^m(1)}\nonumber\\
&&\qquad +\int_0^tK^m(t-s)f^m(s)ds\cdot \beta^m(0),\nonumber\\
\mathcal{Y}^m(t)
&=&\exp\left(-\frac{\gamma t}{m}\right)\mathcal{Y}^m(0)+
\int_{0}^{t}\exp\left(-\frac{\gamma (t-s)}{m}\right)\label{4.27}\\
&&\qquad\times \left(\left(u_{X}^m(s)+\beta^m(0)\right)f^m(s)ds +dW^m(s)\right)\nonumber\\
&=&Y^m(t)-\exp\left(-\frac{\gamma t}{m}\right)\frac{X(\varphi^m(0))-X(0)-K^m(1)\mathcal{Y}^m(0)}{K^m(1)}\nonumber\\
&&\qquad +\int_{0}^{t}\exp\left(-\frac{\gamma (t-s)}{m}\right)f^m(s)ds\cdot \beta^m(0).\nonumber
\end{eqnarray}
From (\ref{2.29.00}) and (\ref{4.35}), 
\begin{equation}
\mathcal{X}^m(1)=X^m(1)=X(1),\quad \mathcal{Z}^m=(\mathcal{X}^m, \mathcal{Y}^m)\in \mathcal{A}^m(P_0,P_1)
\end{equation}
(see (\ref{2.12}), (\ref{3.27}), and (\ref{2.44}) for notation).
Indeed,
$$\int_0^1K^m(1-s)f^m(s)ds=\frac{1-\varphi^m(0)}{\gamma}.$$

From Theorem \ref{thm2.3}, Lemma \ref{lm3.5}, and (\ref{4.35})--(\ref{4.27}),
$(\mathcal{X}^m,\mathcal{Y}^m)$ converges  to $(X,0)$, as $m\to 0$, locally uniformly on $[0,1)$, a.s..
Indeed,  as $m\to0$,
\begin{eqnarray}
\quad \|\mathcal{X}^m-X^m\|_\infty&\le& 2(|X(\varphi^m(0))-X(0)|+
\gamma^{-1}|\mathcal{Y}^m(0)|)\to 0, \\
\|\mathcal{Y}^m-Y^m\|_\infty&\le& \left(\frac{1}{K^m(1)}+\frac{\gamma}{1-\varphi ^m(0)}\right)
 (|X(\varphi^m(0))-X(0)|+\gamma^{-1}|\mathcal{Y}^m(0)|)\nonumber\\
 \\
 &\to& 0, \nonumber
 \end{eqnarray}
since 
$$0\le K^m(t)\le K^m(1)<\frac{1}{\gamma}, \quad \gamma K^m(t-s)\le f^m(s)\le 1.$$

%We prove (\ref{2.29}).
From (A2, iii),
\begin{eqnarray}\label{4.35.1}
&&E\left[\int_0^1L\left(t, \mathcal{Z}^m(t);\left(u_{X}^m(t)+\beta^m(0)\right)f^m(t)\right)dt\right]\\
& \le& 
E\left[\int_0^1R_1\left(t, \mathcal{Z}^m(t); u_X(\varphi^m(t))+\beta^m(0)\right)f^m(t)^2dt\right]
+E\left[\int_0^1L(t, \mathcal{Z}^m(t); 0)dt\right]\nonumber\\
&=&
E\left[\int_{0}^1L\left(t, \mathcal{Z}^m(t); u_X(\varphi^m(t))+\beta^m(0)\right)f^m(t)^2dt\right]\nonumber\\
&&\qquad +E\left[\int_0^1(-f^m(t)^2+1)L(t, \mathcal{Z}^m(t); 0)dt\right].\nonumber
\end{eqnarray}
From (A4),
\begin{eqnarray}
&&L\left(t, \mathcal{Z}^m(t); u_X(\varphi^m(t))+\beta^m(0)\right)\\
&\le&L\left(t, \mathcal{Z}^m(t); u_X(\varphi^m(t))\right)
+C|\beta^m(0)|\left(|u_X(\varphi^m(t))|^{r_0-1}
+\left|\beta^m(0)\right|^{r_0-1}\right).\nonumber
\end{eqnarray}
By H\"older's inequality,
\begin{eqnarray}%\label{4.38}
&&E\left[|\beta^m(0)| \int_0^1|u_X(\varphi^m(t))|^{r_0-1}f^m(t)^2dt\right]\\
&=&E\left[|\beta^m(0)| \int_{\varphi^m(0)}^1|u_X(s)|^{r_0-1}ds\right]\nonumber\\
&\le&\{E [|\beta^m(0)|^{r_0}]\}^{1/r_0}
\left \{E\left[\int_{\varphi^m(0)}^1|u_X(s)|^{r_0}ds\right]\right\}^{(r_0-1)/r_0}.\nonumber
\end{eqnarray}
\begin{eqnarray}%\label{4.43}
&&
\left|1-\varphi ^m(0)\right|^{r_0}
E\left[\left|\beta^m(0)\right|^{r_0}\right]\\
&\le & 3^{r_0-1}
E\left[\left(\int_0^{\varphi^m(0)}|u_X(s)|ds\right)^{r_0}+|W(\varphi^m(0))|^{r_0}
+|\mathcal{Y}^m(0)|^{r_0}\right]\nonumber\\
&\le & 3^{r_0-1}
E\left[\varphi^m(0)^{r_0-1}\int_0^{\varphi^m(0)}|u_X(s)|^{r_0}ds
+\varphi^m(0)^{r_0/2}|W(1)|^{r_0}
+|\mathcal{Y}^m(0)|^{r_0}\right].\nonumber
\end{eqnarray}
For sufficiently large $R>0$,
\begin{eqnarray}
&&E\left[\int_0^{\varphi^m(0)}|u_X(s)|^{r_0}ds\right]\\
&\le & \left(\varphi^m(0)R^{r_0}+\frac{1}{C_{r_0,R}}E\left[\int_{0}^{\varphi^m(0)}L_0(t, X(t);u_X(t))dt\right]\right) \to0, \quad m\to 0,\nonumber
\end{eqnarray}
\begin{equation}\label{4.40}
E\left[\int_{\varphi^m(0)}^1|u_X(s)|^{r_0}ds\right]
\le  R^{r_0}+\frac{1}{C_{r_0, R}}E\left[\int_{0}^1L_0(t, X(t);u_X(t))dt\right]<\infty.
\end{equation}

From (\ref{4.35.1})--(\ref{4.40}),  (\ref{2.48})--(\ref{2.49}) can be proven in the same way as Theorem \ref{thm2.3}.
\end{proof}



From Lemma \ref{lm3.3}, in the same way as Theorem \ref{pp1.1},
we can prove Proposition \ref{pp2.2}.

\begin{proof}[Proof of Proposition \ref{pp2.2}]
From Theorem 2.2.15 and Lemma 3.2.3 in \cite{DS}, the convex duality does the proof.
We only have to prove that $P\mapsto V^m(B,P_0;P)$ is convex and lower semicontinuous.

We first prove that $P\mapsto V^m(B,P_0;P)$ is convex.
For any  $Z_n=(X_n, Y_n)\in \mathcal{A}^m(B, P_0;P^{X_n(1)})$ such that $E[\int_0^1 |u_{X_n} (t)|dt]$ are finite, $n=0,1$,  and  for any $\lambda\in (0,1)$, let 
$$Q_\lambda(dt\hbox{ }dz\hbox{ }du):=(1-\lambda) dt P^{(Z_0(t),u_{X_0} (t))}(dz\hbox{ }du)
 +\lambda dt P^{(Z_1(t),u_{X_1} (t))}(dz\hbox{ }du).
$$
Then, in the same way as the proof of Theorem \ref{pp1.1},
 there exists $Z_\lambda \in \mathcal{A}^m$ such that 
\begin{eqnarray}
%dt P^{(Z_\lambda(t),u_{X_\lambda} (t))}(dz\hbox{ }du)=&(1-\lambda) dt P^{(Z_0(t),u_{X_0} (t))}(dz\hbox{ }du)\\
%&&\qquad +\lambda dt P^{(Z_1(t),u_{X_1} (t))}(dz\hbox{ }du),\\
u_{X_\lambda} (t)&=&
E^{Q_\lambda}[u|t,Z_\lambda(t)],\\
P^{Z_\lambda(t)}(dz)&=&(1-\lambda) P^{Z_0(t)}(dz)+\lambda P^{Z_1(t)}(dz),
\quad 0\le t\le 1
\end{eqnarray}
by Bogachev--R\"ockner--Shaposhnikov's superposition principle (see \cite{Bog21}, Theorem 1).
Here
$E^{Q_\lambda}[u|t,z]$ denotes the conditional expectation of $u$ given $(t,z)$ under $Q_\lambda$.
Indeed, replacing $P^{Z_\infty (t)}$ and $Q_\infty$ by $(1-\lambda) P^{Z_0(t)}+\lambda P^{Z_1(t)}$
and $Q_\lambda$, respectively, (\ref{4.3.0}) holds by the It\^o formula.
Since $B$ is convex, $Z_\lambda\in \mathcal{A}^m(B, P_0;(1-\lambda) P^{X_0(1)}
+\lambda  P^{X_1(1)})$.
In particular, the following holds and implies that $P\mapsto V^m(B,P_0;P)$ is convex:
\begin{eqnarray}
&&(1-\lambda) E\biggl[\int_0^1 L(t,Z_0(t);u_{X_0} (t))dt \biggr]
+\lambda E\biggl[\int_0^1 L(t,Z_1(t);u_{X_1} (t))dt \biggr]\\
&=&\int_{[0,1]\times \mathbb{R}^{2d}\times \mathbb{R}^{d} }
L(t,z;u)Q_\lambda(dt\hbox{ }dz\hbox{ }du)\nonumber\\
&\ge& \int_{[0,1]\times \mathbb{R}^{2d}}
L(t,z;E^{Q_\lambda}[u|t,z])Q_\lambda(dt\hbox{ }dz\times \mathbb{R}^{d} )
=E\biggl[\int_0^1 L(t,Z_\lambda(t);u_{X_\lambda} (t))dt \biggr],\nonumber
\end{eqnarray}
by Jensen's inequality.
From Lemma \ref{lm3.3}, in the same way as Theorem \ref{pp1.1},
one can also show the lower--semicontinuity of $P\mapsto V^m(B,P_0;P)$.
\end{proof}

For $Z=(X,Y)\in \mathcal{A}^m$,
the SDE for $X(t)+K^m(1-t)Y(t)$ plays a crucial role in the proof of Proposition \ref{pp2.3}.

\begin{proof}[Proof of Proposition \ref{pp2.3}]
For $Z=(X,Y)\in \mathcal{A}^m(P_0,P^{X(1)})$, let
\begin{equation}\label{2.34}
\eta(t):=X(t)+K^m(1-t)Y(t)
\end{equation}
(see (\ref{3.7}) for notation).
Then
\begin{eqnarray}
d\eta (t)&=&K^m(1-t)(u_X(t)dt+dW(t)),\label{2.35}\\
%\eta(0)&=&X(0)+K^m(1)Y(0),\nonumber\\
\eta(1)&=&X(1).\nonumber
\end{eqnarray}
Indeed, by the It\^o formula, for $t\in (0,1)$,
\begin{eqnarray*}
d\{X(t)+K^m(1-t)Y(t)\}
&=&\frac{1}{m}Y(t)dt-\frac{1}{m}\exp\left(-\frac{\gamma}{m}(1-t)\right)Y(t)dt\\
&&\qquad +K^m(1-t)\left\{\left(u_X(t)-\frac{\gamma }{m}Y(t)\right)dt+dW(t)\right\}.
\end{eqnarray*}
From (\ref{2.35}), $u_\eta =u_X$, and 
\begin{equation}
E\biggl[f(X (1))-\int_0^1 L(t;u_X(t))dt \biggr]
=E\biggl[f(\eta (1))-\int_0^1 L(t;u_\eta(t))dt \biggr],
\end{equation}
which implies that the l.h.s. is less than or equal to the r.h.s. in (\ref{2.34.1}).

Suppose that $\{\eta (t)\}_{0\le t\le 1}$ be a semimartingale defined on a complete filtered probability space such that the following holds:
\begin{equation}
d\eta (t)=K^m(1-t)\left\{u_\eta(t)dt+dW(t)\right\}.%\label{2.35}\\
%P^{\eta(1)}&=&P_1.\nonumber
\end{equation}
Let $X_0$  be an $\mathbb{R}^d$--valued random variable defined
on the same probability space as  $\{\eta (t)\}_{0\le t\le 1}$ such that $P^{X_0}=P_0$.
Let $Z=(X,Y)$ be a solution to (\ref{1.1})--(\ref{1.2}) such that 
$$
u_X(t):=u_\eta(t),\quad X(0)=X_0, \quad Y(0)=\frac{\eta(0)-X_0}{K^m(1)}.$$
Then
\begin{eqnarray}
d\eta (t)&=&K^m(1-t)\left\{u_X(t)dt+dW(t)\right\},\label{4.37}\\
%\eta(0)&=&X(0)+K^m(1)Y(0),\nonumber\\
\eta(1)&=&X(1).\nonumber
\end{eqnarray}
Indeed, from (\ref{2.34})--(\ref{2.35}), 
\begin{eqnarray*}
X(1)&=&X(1)+K^m(1-1)Y(1)\\
&=&X(0)+K^m(1-0)Y(0)
+\int_0^1 d(X(t)+K^m(1-t)Y(t))\\
&=&\eta(0)+\int_0^1 K^m(1-t)\left\{u_X(t)dt+dW(t)\right\}\\
&=&\eta(1).
\end{eqnarray*}
From (\ref{4.37}),
\begin{equation}
E\biggl[f(\eta (1))-\int_0^1 L(t;u_\eta(t))dt \biggr]
=E\biggl[f(X (1))-\int_0^1 L(t;u_X(t))dt \biggr].
\end{equation}
Since $P^{X(0)}=P^{X_0}=P_0$,
the l.h.s. is greater than or equal to the r.h.s. in  (\ref{2.34.1}).

If $L=|u|^2/2$, then $V^0(P_0,P_1)$ is Schr\"odinger's problem  (see section \ref{sec:1}).
If $V^m(P_0,P_1)$  converges to $V^0(P_0,P_1)$ as $m\to 0$,
then $V^0(P_0,P_1)$ does not depend on $P_0$ either, which is a contradiction
 (see \cite{M2021} for more general SOTs).
 
\end{proof}


%\newpage
\begin{thebibliography}{99.}%

\bibitem{Birrell}
\textsc{Birrell, J.} and \textsc{Wehr, J.} (2020).
Langevin equations in the small-mass limit: higher-order approximations.
\textit{Ann. Henri Poincar\'e} 
\textbf{21, no. 6} 1765--1811.

\bibitem{Bog21}
\textsc{Bogachev, V. I.}, \textsc{R\"ockner, M.} and \textsc{Shaposhnikov, S. V.}
 (2021).
On the Ambrosio--Figalli--Trevisan superposition principle for probability solutions to Fokker--Planck--Kolmogorov equations.
\textit{J. Dynam. Differential Equations} \textbf{33} 715--739.


\bibitem{Bre1}
\textsc{Brenier, Y.} (1987).
D\'ecomposition polaire et r\'earrangement monotone des champs de vecteurs. 
\textit{C. R. Acad. Sci. Paris S\'erie I } \textbf{305, no. 19} 805--808.

\bibitem{Bre2}
\textsc{Brenier, Y.} (1991).
Polar factorization and monotone rearrangement of vector-valued functions.
\textit{Comm. Pure Appl. Math.} \textbf{44,  no. 4} 375--417.

\bibitem{Cover}
\textsc{Cover, T. M.} and \textsc{Thomas, J. A.} (2006).
\textit{Elements of Information Theory}, 2nd ed. 
Wiley--Interscience, New York.

\bibitem{D91}
\textsc{Dai Pra, P.} (1991).
A stochastic control approach to reciprocal diffusion processes.
\textit{Appl. Math. Optim.}  \textbf{23} 313--329.

\bibitem{DS}
\textsc{Deuschel, J. D.} and \textsc{Stroock} D. W. (1989).
\textit{Large Deviations}.
Academic Press, Inc., Boston, MA.

%\bibitem{FS06}
%\textsc{Fleming, W. H.} and \textsc{Soner, H. M.} (2006).
%\textit{Controlled Markov Processes and Viscosity Solutions}, 2nd ed. 
%Springer--Verlag, New York. 

\bibitem{F88} 
\textsc{F\"ollmer, H.} (1988).
Random fields and diffusion processes.
\textit{In: Hennequin, P.L. (ed.)
\'Ecole d'\'Et\'e de Probabilit\'es de Saint-Flour XV--XVII, 1985--87,  Lecture Notes in Math.}  \textbf{1362} 101--203.
Springer--Verlag, Berlin.

\bibitem{FW-1} 
\textsc{Freidlin, M.} and \textsc{Hu, W.} (2011).
Smoluchowski-Kramers approximation in the case of variable friction.
\textit{J. Math. Sci.} \textbf{179} 184--207.

\bibitem{FW-2} 
\textsc{Freidlin, M.}, \textsc{Hu, W.} and \textsc{Wentzell, A.} (2013).
Small mass asymptotic for the motion with vanishing friction.
\textit{Stochastic Process. Appl.} \textbf{123, no. 1} 45--75.


\bibitem{GangboMcCann}
\textsc{Gangbo, W.} and \textsc{McCann, R. J.} (1996).
The geometry of optimal transportation. 
\textit{Acta Math.} \textbf{177, no. 2} 113--161.

\bibitem{IW14}
\textsc{Ikeda, N.} and \textsc{Watanabe, S.} (2014).
\textit{Stochastic Differential Equations and Diffusion Processes}, 2nd ed.
North--Holland/Kodansha, Tokyo.

\bibitem{Ishii}
\textsc{Ishii, H.}, \textsc{Souganidis, P. E.} and \textsc{Tran, H. V.} (2017).
 On the Langevin equation with variable friction. 
\textit{Calc. Var.} \textbf{56} 161.

\bibitem{2-JS}
\textsc{Jacod, J.} and \textsc{Shiryaev, A. N.} (1987).
\textit{Limit Theorems for Stochastic Processes}.
Springer--Verlag, Heidelberg.

\bibitem{1-Jamison1974}
\textsc{Jamison, B.} (1974).
Reciprocal processes. 
\textit{Z. Wahrscheinlichkeitstheorie und Verw. Gebiete} \textbf{30} 65--86.

\bibitem{J75}
\textsc{Jamison, B.} (1975).
The Markov process of Schr\"odinger.
\textit{Z. Wahrscheinlichkeitstheorie und Verw. Gebiete} \textbf{32} 323--331.

\bibitem{L2} 
\textsc{L\'eonard, C.}  (2014).
A survey of the Schr\"odinger problem and some of its connections with optimal transport.
Special Issue on Optimal Transport and Applications. 
\textit{Discrete Contin. Dyn. Syst.} \textbf{34} 1533--1574.

%\bibitem{L}
%L\"ocherbach, E.:
%Ergodicity and speed of convergence to equilibrium for diffusion processes,
%https://eloecherbach.u-cergy.fr/cours.pdf, Proposition 4.12, November 14, 2015.

\bibitem{M04}
\textsc{Mikami, T.} (2004).
Monge's problem with a quadratic cost by the zero-noise limit of h-path  processes.
\textit{Probab. Theory Related Fields} \textbf{129} 245--260.

\bibitem{M08}
\textsc{Mikami, T.}  (2008).
Marginal problem for semimartingales via duality.
\textit{In: Giga, Y., Ishii, K., Koike, S., Ozawa, T., Yamada, N. (eds.)
International Conference for the 25th Anniversary of Viscosity Solutions.
Gakuto International Series. Mathematical Sciences and Applications}
\textbf{30} 133--152.
Gakkotosho, Tokyo.

\bibitem{M09} 
\textsc{Mikami, T.}  (2009).
Optimal transportation problem as stochastic mechanics.
\textit{In Selected Papers on Probability and Statistics, Amer. Math. Soc. Transl.} \textbf{Ser. 2, 227} 75--94.
 Amer. Math. Soc., Providence.
 
\bibitem{M21}
\textsc{Mikami, T.}  (2021).
Stochastic optimal transport revisited.
\textit{SN Partial Differ. Equ. Appl.} \textbf{2} article 5, 26 pages.

\bibitem{M2021}
\textsc{Mikami, T.} (2021).
\textit{Stochastic Optimal Transportation: Stochastic Control with Fixed Marginals}.
SpringerBriefs in Mathematics, Springer, Singapore.

\bibitem{MT06}
\textsc{Mikami, T.} and \textsc{Thieullen, M.}  (2006).
Duality theorem for stochastic optimal control problem.
\textit{Stochastic Process. Appl.} \textbf{113} 1815--1835.


\bibitem{1-MTsiam}
\textsc{Mikami, T.} and \textsc{Thieullen, M.}  (2008).
Optimal transportation problem by stochastic optimal control.
\textit{SIAM J. Control Optim.} \textbf{47, no. 3} 1127--1139.
 

\bibitem{PAV}
\textsc{Pavliotis, G. A.} (2014).
\textit{Stochastic Processes and Applications: Diffusion Processes, the Fokker-Planck and Langevin Equations}.
Springer Science+Business Media, New York.

\bibitem{1-PC}
\textsc{Peyr\'e, G.} and \textsc{Cuturi, M.} (2019)
\textit{Computational Optimal Transport: With Applications to Data Science}.
Now Publishers, Boston.

\bibitem{1-RR}
\textsc{Rachev, S. T.} and \textsc{R\"uschendorf, L.}  (1998).
\textit{Mass Transportation Problems, Vol. I: Theory, Vol. II: Application}.
Springer--Verlag, Heidelberg.

\bibitem{1-RT}
\textsc{R\"uschendorf, L.} and \textsc{Thomsen, W.} (1993).
Note on the Schr\"odinger equation and $I$-projections.
\textit{Statist. Probab. Lett.} \textbf{17, no. 5} 369--375.

\bibitem{S1}
\textsc{Schr\"odinger, E.} (1931).
\"Uber die Umkehrung der Naturgesetze.
\textit{Sitz. Ber. der Preuss. Akad. Wissen., Berlin, Phys. Math.}  \textbf{8, no. 9} 144--153.

\bibitem{S2}
\textsc{Schr\"odinger, E.} (1932).
Th\'eorie relativiste de l'electron et l'interpr\'etation de la m\'ecanique quantique.
\textit{Ann. Inst. H. Poincar\'e} \textbf{2} 269--310 .

\bibitem{1-V2}
\textsc{Villani, C.}  (2008).
\textit{Optimal Transport: Old and New}.
Springer--Verlag, Heidelberg.

\bibitem{Z86}
\textsc{Zambrini, J. C.} (1986).
Variational processes.
\textit{In:  Albeverio, S., Casati, G., Merlini, D. (eds.)
Stochastic Processes in Classical and Quantum Systems, Ascona 1985,
Lecture Notes in Phys.} \textbf{262} 517--529.
Springer--Verlag, Heidelberg.

\bibitem{Z86-2}
\textsc{Zambrini, J. C.} (1986).
Variational processes and stochastic versions of mechanics.
\textit{J. Math. Phys.} \textbf{27} 2307--2330.

\bibitem{Z86-3}
\textsc{Zambrini, J. C.} (1986).
Stochastic mechanics according to E. Schr\"odinger.
\textit{Phys. Rev. A(3)} \textbf{33} 1532--1548.

\end{thebibliography}


\end{document}
